% Generated by code/build_integration_fragment.py; edit the modular sources.
% Research continuation, 10 October 2026; all labels have prefix mld:.
% This is a BODY fragment, not a standalone document.
% Standard dependencies: amsmath, amssymb, amsthm, mathtools, tabularx,
% booktabs, array, enumitem, xurl, hyperref. Define theorem, proposition,
% corollary, lemma, example and remark environments in the host preamble.
% Copy or reconcile the notation definitions in the standalone main preamble:
% Li, FP, CT, Res, dd, ii, C, R, Q, Z, N, T, e, sgn, lcm, src, zetaj, Lij.
% Column Y is a ragged-right X column. Bibliography entries: references.tex.
% Compilation inside the canonical manuscript is an integration step.

% SOURCE: sections/01-scope.tex
\section{Scope, source audit, and mathematical status}\label{mld:sec:scope}

The canonical manuscript already contains extensive polylogarithmic reduction, gamma arithmetic, Hurwitz-parameter differentiation, and Stieltjes antiderivative theory. The latest incoming reports add periodic Hurwitz convolution, shifted correlations, their spectral derivatives, and contact terms under differentiation. A useful continuation must therefore distinguish an additional identity from a repackaging of an existing result.

We inspected the manuscript text and status files and all five archives in the incoming directory at commit
\src{fc4d3bf80534ad7c901d3b8c9e71baf2df0064ed}
\cite{ProveIt,IncomingResonance,IncomingShifted,IncomingConvolution,IncomingClosure,IncomingCalculus}.
The source manifest records the retrieved bytes and their hashes. The audit was directed at the identity questions relevant to this article; it is not a claim to have re-proved every theorem in the manuscript.

\begin{center}
\begin{tabularx}{\textwidth}{@{}>{\raggedright\arraybackslash}p{0.24\textwidth}Y@{}}
\toprule
Source strand & Consequence for the present work\\
\midrule
Manuscript integration and differentiation chapters & The elementary-symmetric harmonic derivative tower and Hurwitz-jet antiderivatives are established starting points.\\
Manuscript Gaussian harmonic sums & The recorded $S_4$ evaluation has an exact proof. The residual certificates for $S_6$ and $S_8$ do not prove equality.\\
Incoming shifted Hurwitz jets & The explicit question about dilation, traces, and subtraction scale is answered in Section~\ref{mld:sec:dilation}.\\
Other four incoming packages & Their convolution, resonance, and correlation identities constrain the novelty and normalization claims here.\\
Published Euler-sum literature & The unshifted parity reductions in Section~\ref{mld:sec:harmonic} are consistent with Xu--Wang's prior formulas.\\
\bottomrule
\end{tabularx}
\end{center}

\subsection{What is proved here}

The main additions are an explicit finite transform formula for arbitrary denominator order, its complete treatment at integer Mellin resonances, a shifted harmonic reflection generator, and a dilation theorem that retains the local finite-part terms. These are identity results. Estimates appear only when needed to justify convergence, parameter differentiation, or continuation.

The Mellin formula is derived from the classical Fermi integral, Euler's beta integral, and elementary partial fractions. The higher-pole reduction and its integer-parameter formulas organize those ingredients into a uniform calculation. The two-shift harmonic theorem packages an infinite family into one identity of analytic functions. The dilation result closes a question explicitly proposed in the incoming material, with a precise choice of subtraction coordinate.

None of these descriptions asserts global literature priority for every specialization. The standard special-function conventions follow the DLMF \cite{DLMFpoly,DLMFlerch,DLMFhurwitz,DLMFgamma}. The relevant Euler-sum antecedent is Xu and Wang \cite{XuWang}; the underlying Hurwitz product-integral method goes back at least to the work of Espinosa and Moll \cite{EspinosaMoll}. The article supplies its own proofs of every formula used as a new addition to the repository.

\subsection{Conventions}

For $|z|<1$ and $\Re q>0$,
\[
 \Li_s(z)=\sum_{n\ge1}\frac{z^n}{n^s},\qquad
 \Phi(z,s,q)=\sum_{n\ge0}\frac{z^n}{(n+q)^s}.
\]
Thus $\Li_s(z)=z\Phi(z,s,1)$. Powers of $n+q$ use the principal logarithm on the right half-plane; more generally principal logarithms are used unless a continuation is specified. We write $\psi(q)=\Gamma'(q)/\Gamma(q)$. The gamma logarithm is the analytic $\log\Gamma(q)$ on $\Re q>0$, normalized by $\log\Gamma(1)=0$.

Spectral derivatives are denoted by square brackets:
\[
 \Li_s^{[m]}(z)=\partial_s^m\Li_s(z),\qquad
 \zeta^{[m]}(s,q)=\partial_s^m\zeta(s,q).
\]
An ordinary superscript in $\psi^{(r)}(q)$ or $\gamma_m^{(r)}(q)$ means differentiation in $q$. The normalization of the generalized Stieltjes constants is
\begin{equation}\label{mld:eq:stieltjes-definition}
 \zeta(1+u,q)=\frac1u+
       \sum_{m\ge0}\frac{(-1)^m}{m!}\gamma_m(q)u^m,
 \qquad \gamma_0(q)=-\psi(q).
\end{equation}
We write $\gamma_m=\gamma_m(1)$ and $\gamma=\gamma_0$. Generalized harmonic numbers are
\[
 H_n^{(r)}=\sum_{j=1}^n j^{-r},\qquad H_0^{(r)}=0,
 \qquad H_n=H_n^{(1)}.
\]
The rising factorial is $(s)_r=s(s+1)\cdots(s+r-1)$, with $(s)_0=1$. Coefficient extraction, such as $[u^m]F(u)$, always refers to an analytic germ or to a specified formal power series. A removable product such as $u\zeta(1+u)$ is evaluated as a whole.

There are two different regularization operations later in the article. Distributional pullback transports an already extended distribution. Ambient-coordinate finite part instead extends the dilated pointwise function by deleting divergent terms in the original integration coordinate. Their difference is explicit and generally nonzero. Throughout the Mellin and harmonic sections, the integrals are ordinary convergent integrals and no finite-part convention is involved.


% SOURCE: sections/02-mellin.tex
\section{A finite Mellin--Lerch calculus}\label{mld:sec:mellin}

\subsection{The analytic difference that removes the Lerch cut}

Put $\mathcal S=\C\setminus(-\infty,0]$. For $z\in\mathcal S$ we use $\log z$ with imaginary part in $(-\pi,\pi)$. The natural expression in the transform is
\begin{equation}\label{mld:eq:Lerch-difference}
 \mathcal L(s,q;z)=z^q\Phi(z,s,q)-\Li_s(z),\qquad \Re q>0.
\end{equation}
Initially this is defined by the two convergent series for $0<z<1$. The individual summands have the usual positive-axis continuation cut; their difference does not.

\begin{lemma}[Cancellation and the unit-argument germ]\label{mld:lem:Lerch-cancellation}
The expression \eqref{mld:eq:Lerch-difference} extends to a holomorphic function of $z\in\mathcal S$, holomorphic in $\Re q>0$, and entire in $s$. Its value at $z=1$ is
\begin{equation}\label{mld:eq:Lerch-unit}
 \mathcal L(s,q;1)=\zeta(s,q)-\zeta(s),
\end{equation}
where the pole at $s=1$ cancels. For $z>1$, either consistent lateral continuation of the two terms in \eqref{mld:eq:Lerch-difference} gives this same function. Locally at $z=1$, with $\mu=\log z$,
\begin{equation}\label{mld:eq:Lerch-germ}
 \mathcal L(s,q;\e^\mu)
 =\sum_{k\ge0}\frac{\zeta(s-k,q)-\zeta(s-k)}{k!}\mu^k,
 \qquad |\mu|<2\pi,
\end{equation}
understood on the local logarithm sheet. The series and its fixed-order parameter derivatives converge locally uniformly.
\end{lemma}

\begin{proof}
For $\Re s>0$ the Lerch integral writes the first term as
\[
 \frac{z^q}{\Gamma(s)}\int_0^\infty
          \frac{t^{s-1}\e^{-qt}}{1-z\e^{-t}}\dd t.
\]
For $z>1$, the residue at $t=\log z$ is
$(\log z)^{s-1}/\Gamma(s)$, independent of $q$. It is therefore the same as the residue for $q=1$, which gives $\Li_s(z)$. The two lateral jumps cancel. This argument first applies away from exceptional spectral orders and extends to them analytically.

At $\mu=0$ the classical Lerch expansion \cite[\S25.14]{DLMFlerch} has the form
\[
 \e^{q\mu}\Phi(\e^\mu,s,q)
 =\Gamma(1-s)(-\mu)^{s-1}
       +\sum_{k\ge0}\zeta(s-k,q)\frac{\mu^k}{k!}.
\]
Subtract its $q=1$ version. The nonanalytic term is independent of $q$ and cancels exactly, giving \eqref{mld:eq:Lerch-germ}. At integer $s$, subtract first and continue in $s$; each coefficient is a difference of two Hurwitz zeta functions with the same residue. This proves regularity at $z=1$ and \eqref{mld:eq:Lerch-unit}. Local uniform convergence follows from the convergent Lerch expansion on smaller $\mu$-disks and Cauchy estimates on compact parameter sets. The series definition gives entire dependence on $s$ for $0<z<1$, and continuation preserves it. There are no remaining obstructions in $\mathcal S$: the positive-axis jumps have cancelled, $z=1$ is removable, and $z=0$ is excluded by the chosen domain. Equivalently, the integral representation in Theorem~\ref{mld:thm:Mellin-master} below independently supplies this continuation when $0<\Re q<2$, after which the finite Lerch shift identity extends it to every $\Re q>0$.
\end{proof}

\begin{remark}
The local logarithm qualification in \eqref{mld:eq:Lerch-germ} matters. It is a power series in $\mu$, not a license to replace $\log z$ by an arbitrary logarithm while keeping the other branches fixed. In computations at $z>1$, taking inconsistent lateral values before subtraction introduces a spurious imaginary term.
\end{remark}

\subsection{The master transform at every denominator order}

For an integer $N\ge1$ define
\begin{equation}\label{mld:eq:JN-definition}
 J_N(a,s;z)=\int_0^\infty
       \frac{x^{a-1}\Li_s(-zx)}{(1+x)^N}\dd x.
\end{equation}
Let the rational polynomials $q_{N,j}(a)$ be determined by
\begin{equation}\label{mld:eq:q-polynomial}
 Q_{N,a}(X):=\binom{X+a-1}{N-1}
 =\frac1{(N-1)!}\prod_{r=1}^{N-1}(X+a-r)
 =\sum_{j=0}^{N-1}q_{N,j}(a)X^j.
\end{equation}
For $N=1$ the product is $1$.

\begin{theorem}[Finite Mellin--Lerch transform]\label{mld:thm:Mellin-master}
Let $N\ge1$, $z\in\mathcal S$, $s\in\C$, and
\begin{equation}\label{mld:eq:Mellin-strip}
 -1<\Re a<N.
\end{equation}
Then \eqref{mld:eq:JN-definition} is an ordinary convergent integral, holomorphic in these parameters and entire in $s$. For noninteger $a$,
\begin{equation}\label{mld:eq:Mellin-master}
 \boxed{\displaystyle
 J_N(a,s;z)=\frac{(-1)^N\pi}{\sin\pi a}
       \sum_{j=0}^{N-1}q_{N,j}(a)
          \mathcal L(s-j,N-a;z).}
\end{equation}
All the apparent singularities at $a=0,1,\ldots,N-1$ are removable. The same formula, continued as a whole, holds there. The strip \eqref{mld:eq:Mellin-strip} is the maximal common strip for all spectral orders: already $s=0,z=1$ forces both endpoints.
\end{theorem}

\begin{proof}
We first establish uniform convergence and spectral holomorphy, including negative spectral orders. Fix compact sets of $s$ and $z\in\mathcal S$. Choose an integer $r\ge0$ such that $\Re(s+r)>0$ throughout a neighborhood of the spectral compact set. The Fermi integral gives
\[
 \Li_{s+r}(-zx)=-\frac1{\Gamma(s+r)}
   \int_0^\infty t^{s+r-1}\frac{zx\e^{-t}}{1+zx\e^{-t}}\dd t.
\]
Apply $(x\partial_x)^r$ to obtain $\Li_s(-zx)$. The differentiated rational kernel is bounded by a constant times $\min(x\e^{-t},1)$, uniformly in these compact sets: the ray of $zx\e^{-t}$ remains a positive angular distance from the negative real axis, and its denominator is bounded away from zero relative to $1+|zx\e^{-t}|$. Integrating this bound, splitting at $t=\log x$ when $x>1$, shows
\[
 \Li_s(-zx)=O(x)\quad(x\downarrow0),\qquad
 \Li_s(-zx)=O((1+\log x)^M)\quad(x\to\infty)
\]
for some finite $M$, uniformly on the chosen compacts. These estimates give an integrable majorant in every smaller strip in \eqref{mld:eq:Mellin-strip}. They justify holomorphy of \eqref{mld:eq:JN-definition}. They also justify arbitrary fixed powers of $\log x$, and spectral differentiation follows either directly or from Cauchy estimates. For $s=0,z=1$, $\Li_0(-x)=-x/(1+x)$ shows the sharpness of the common strip.

For $N=1$, take initially $0<z<1$, $0<\Re a<1$, and $\Re s>0$. Fubini's theorem applies to the Fermi integral. If $b>0$, partial fractions and the beta integral give
\begin{equation}\label{mld:eq:beta-resolvent}
 \int_0^\infty\frac{x^a}{(1+x)(x+b)}\dd x
       =\frac{\pi}{\sin\pi a}\frac{b^a-1}{b-1}.
\end{equation}
At $b=1$ the right side is interpreted by continuity. One can prove \eqref{mld:eq:beta-resolvent} first for $-1<\Re a<0$, where the two partial fractions are separately integrable, and then continue to $-1<\Re a<1$, where their difference is integrable. This avoids separating divergent integrals.

Substitute $b=\e^t/z$ into \eqref{mld:eq:beta-resolvent}. The result is
\begin{align*}
 J_1(a,s;z)
 &=-\frac{\pi}{\Gamma(s)\sin\pi a}
   \int_0^\infty t^{s-1}
      \frac{z^{1-a}\e^{at}-z}{\e^t-z}\dd t\\
 &=-\frac{\pi}{\sin\pi a}
       \big[z^{1-a}\Phi(z,s,1-a)-\Li_s(z)\big].
\end{align*}
This proves \eqref{mld:eq:Mellin-master} for $N=1$ on an open set.

Let $E$ be the spectral backward shift, $(Ef)(s)=f(s-1)$. Integration of the derivative of $x^a(1+x)^{-N}\Li_s(-zx)$, with vanishing endpoint terms, gives
\begin{equation}\label{mld:eq:JN-recurrence}
 J_{N+1}(a,s;z)
       =\frac{(N-a)J_N(a,s;z)-J_N(a,s-1;z)}{N}.
\end{equation}
Consequently, in a common strip,
\[
 J_N=\frac1{(N-1)!}\prod_{r=1}^{N-1}(r-a-E)J_1.
\]
Applying this polynomial in $E$ to the Lerch series annihilates its first $N-1$ terms, because
\begin{equation}\label{mld:eq:annihilation}
 Q_{N,a}(n+1-a)=\binom n{N-1}=0
       \quad(0\le n<N-1).
\end{equation}
Shift the remaining index by $N-1$. The Lerch parameter becomes $N-a$, and the power $z^{1-a}$ becomes $z^{N-a}$. Since
\[
 Q_{N,a}(n+N-a)=\binom{n+N-1}{N-1},
\]
the resulting expression is exactly \eqref{mld:eq:Mellin-master}. The parameter $N-a$ has positive real part on the entire strip \eqref{mld:eq:Mellin-strip}. Holomorphic continuation in $a$ and $s$, followed by Lemma~\ref{mld:lem:Lerch-cancellation} in $z$, proves the identity throughout the asserted domain. The original integral is holomorphic at every integer in that strip, so the remaining cosecant poles are removable. The next theorem gives their values without relying on an unspecified limit.
\end{proof}

The first two cases illustrate the finite character of the result:
\begin{align}
 J_1(a,s;z)&=-\pi\csc(\pi a)\mathcal L(s,1-a;z),\label{mld:eq:J1}\\
 J_2(a,s;z)&=\pi\csc(\pi a)
       \big[\mathcal L(s-1,2-a;z)
                    +(a-1)\mathcal L(s,2-a;z)\big].\label{mld:eq:J2}
\end{align}
For general $N$, only $N$ Lerch differences occur. The shift to $N-a$ is essential: retaining $1-a$ term by term would leave the safe Hurwitz half-plane when $\Re a>1$ and obscure the cancellations in \eqref{mld:eq:annihilation}.

\subsection{Integer Mellin parameters and complete logarithmic moments}

Define the entire spectral function
\begin{equation}\label{mld:eq:F-definition}
 F_s(z)=s\Li_{s+1}(z)-\log z\,\Li_s(z),
\end{equation}
with the same consistent continuation across $z>1$. At $z=1$ this means
\begin{equation}\label{mld:eq:F-unit}
 F_s(1)=s\zeta(s+1),\qquad F_0(1)=1.
\end{equation}
The last value is a removable value, not the undefined product $0\cdot\zeta(1)$.

\begin{theorem}[Every integer resonance]\label{mld:thm:integer-resonance}
For $0\le m\le N-1$,
\begin{equation}\label{mld:eq:integer-resonance}
 \boxed{\displaystyle
 J_N(m,s;z)=(-1)^{N+m}
       \sum_{j=0}^{N-1}q_{N,j}(m)F_{s-j}(z).}
\end{equation}
More strongly, for sufficiently small $\delta$,
\begin{equation}\label{mld:eq:local-resonance}
 J_N(m+\delta,s;z)=(-1)^{N+m-1}
       \sum_{j=0}^{N-1}q_{N,j}(m+\delta)
                                  J_1(\delta,s-j;z).
\end{equation}
Thus every $\log^h x$ moment at an integer Mellin parameter is a finite combination of ordinary polylogarithms, powers of $\log z$, and even zeta values, with coefficients polynomial in $s$ and rational coefficients.
\end{theorem}

\begin{proof}
Write $L=N-m$. The finite Lerch shift formula is
\[
 \mathcal L(s,L-\delta;z)=\mathcal L(s,1-\delta;z)
       -\sum_{k=1}^{L-1}\frac{z^{k-\delta}}{(k-\delta)^s}.
\]
Multiply its shifted spectral versions by $q_{N,j}(m+\delta)$ and sum in $j$. Every term in the finite sum vanishes identically in $\delta$, since
\[
 Q_{N,m+\delta}(k-\delta)
       =\binom{k+m-1}{N-1}=0\quad(1\le k<L).
\]
Using $\sin\pi(m+\delta)=(-1)^m\sin\pi\delta$ proves \eqref{mld:eq:local-resonance}. The first Taylor coefficient in
\[
 \mathcal L(s,1-\delta;z)
 =\e^{-\delta\log z}
       \sum_{r\ge0}\frac{(s)_r}{r!}\Li_{s+r}(z)\delta^r
       -\Li_s(z)
\]
is $F_s(z)$. This series calculation is first made for $0<z<1$,
where every coefficient is finite, and then continued as a combined
analytic germ. At $z=1$ and spectral resonances the removable
products are evaluated only after combination. Hence $J_1(0,s;z)=-F_s(z)$, proving \eqref{mld:eq:integer-resonance}. Further coefficient extraction proves the last assertion, made explicit below.
\end{proof}

Let
\[
 \frac{\pi t}{\sin\pi t}=\sum_{j\ge0}e_{2j}t^{2j},\qquad
 e_0=1,\quad e_{2j}=2(1-2^{1-2j})\zeta(2j)\quad(j\ge1),
\]
and define, for $r\ge1$,
\begin{equation}\label{mld:eq:dr}
 d_r(s,z)=\sum_{v=0}^{r}
       \frac{(-\log z)^{r-v}(s)_v}{(r-v)!v!}\Li_{s+v}(z).
\end{equation}
Then
\begin{equation}\label{mld:eq:M-moments}
 \mathcal M_h(s;z):=\int_0^\infty
       \frac{\log^h x\,\Li_s(-zx)}{x(1+x)}\dd x
 =-h!\sum_{j=0}^{\lfloor h/2\rfloor}e_{2j}\,d_{h+1-2j}(s,z).
\end{equation}
Combining this with \eqref{mld:eq:local-resonance} gives the completely finite rule
\begin{align}\label{mld:eq:all-integer-moments}
 &\int_0^\infty\frac{x^{m-1}\log^h x\,\Li_s(-zx)}{(1+x)^N}\dd x\notag\\
 &\quad=(-1)^{N+m-1}
   \sum_{j=0}^{N-1}\sum_{v=0}^{h}\binom hv
           q_{N,j}^{(v)}(m)\,\mathcal M_{h-v}(s-j;z).
\end{align}
Terms with $v>N-1$ vanish automatically. All parameter differentiations under these integrals are justified by the uniform majorants in the proof of Theorem~\ref{mld:thm:Mellin-master}.

At $z=1$, the especially short version is
\begin{equation}\label{mld:eq:unit-moments}
 \boxed{\displaystyle
 \mathcal M_h(s;1)=-h!\sum_{j=0}^{\lfloor h/2\rfloor}
   e_{2j}\frac{(s)_{h+1-2j}}{(h+1-2j)!}
                      \zeta(s+h+1-2j).}
\end{equation}
For example,
\begin{align*}
 \mathcal M_0(s;1)&=-s\zeta(s+1),\\
 \mathcal M_1(s;1)&=-\frac{s(s+1)}2\zeta(s+2),\\
 \mathcal M_2(s;1)&=-\frac{s(s+1)(s+2)}3\zeta(s+3)
                          -2\zeta(2)s\zeta(s+1),\\
 \mathcal M_3(s;1)&=-\frac{(s)_4}{4}\zeta(s+4)
                          -3\zeta(2)(s)_2\zeta(s+2).
\end{align*}
In particular, the $s=1,h=2$ case gives the ordinary integral
\[
 \int_0^\infty\frac{\log^2x\log(1+x)}{x(1+x)}\dd x=7\zeta(4).
\]

\begin{example}[Two finite higher-pole reductions]
The $N=3$ integer rows are
\begin{align*}
 J_3(1,s;z)&=\tfrac12\big(F_{s-2}(z)-F_{s-1}(z)\big),\\
 J_3(2,s;z)&=-\tfrac12\big(F_{s-2}(z)+F_{s-1}(z)\big).
\end{align*}
At $z=1,s=1$ they equal $-1/4$ and $-3/4$. The necessary input is $F_0(1)=1$; simply dropping a factor proportional to $s-1$ would give wrong answers.
\end{example}

\subsection{Finite algebra, negative orders, and rational kernels}

The coefficient table in \eqref{mld:eq:integer-resonance} is nondegenerate as a finite polynomial transformation.
\begin{proposition}[Invertible integer table]\label{mld:prop:table-invertible}
For $0\le m,j\le N-1$,
\begin{equation}\label{mld:eq:table-determinant}
 \det\big(q_{N,j}(m)\big)
       =\frac{(-1)^{N(N-1)/2}}{\prod_{j=0}^{N-1}j!}.
\end{equation}
Consequently the $N$ Mellin integrals in \eqref{mld:eq:integer-resonance}
determine, by rational linear algebra, the functions
\[
 F_s(z),\quad F_{s-1}(z),\quad\ldots,\quad F_{s-N+1}(z).
\]
\end{proposition}
\begin{proof}
Replace the rows indexed by $m$ by their successive forward differences at $m=0$. The row operation has determinant one. Pascal's identity gives
\[
 \Delta_m^k Q_{N,m}(X)\big|_{m=0}
       =\binom{X-1}{N-1-k}.
\]
These rows have descending degrees $N-1-k$ and leading coefficients $1/(N-1-k)!$. Their coefficient matrix is triangular after reversing the column order. The reversal contributes $(-1)^{N(N-1)/2}$, proving \eqref{mld:eq:table-determinant}.
\end{proof}
This is a statement about a finite coefficient matrix. It makes no claim about arithmetic independence of values at one spectral order.

For a separate exact check at negative spectral orders, let $\left\{\begin{smallmatrix}n\\k\end{smallmatrix}\right\}$ denote a Stirling number of the second kind. The rational expression for the negative-order polylogarithm yields
\begin{equation}\label{mld:eq:negative-beta}
 J_N(a,-r;1)=\sum_{k=0}^{r}(-1)^{k+1}k!
   \left\{\begin{matrix}r+1\\k+1\end{matrix}\right\}
                     B(a+k+1,N-a),\qquad r\ge0.
\end{equation}
Indeed,
\[
 \Li_{-r}(-x)=\sum_{k=0}^{r}(-1)^{k+1}k!
   \left\{\begin{matrix}r+1\\k+1\end{matrix}\right\}
                  \frac{x^{k+1}}{(1+x)^{k+1}},
\]
and each term integrates to the beta factor in \eqref{mld:eq:negative-beta}. Differentiating this finite beta expression gives an independent gamma/polygamma evaluation of the logarithmic moments at those orders. It checks the spectral-pole cancellations without evaluating a singular zeta product numerically.

\begin{corollary}[Proper rational kernels with negative real poles]\label{mld:cor:rational-kernel}
Let
\[
 R(x)=\sum_{\nu=1}^{M}\sum_{k=1}^{N_\nu}
                        \frac{c_{\nu,k}}{(x+b_\nu)^k},\qquad b_\nu>0.
\]
For $-1<\Re a<1$,
\begin{equation}\label{mld:eq:rational-kernel}
 \int_0^\infty x^{a-1}R(x)\Li_s(-zx)\dd x
       =\sum_{\nu,k}c_{\nu,k}b_\nu^{a-k}
                          J_k(a,s;zb_\nu).
\end{equation}
In particular every logarithmic moment at $a=0$ reduces finitely by \eqref{mld:eq:all-integer-moments}; spectral derivatives reduce in the same way. If cancellation in $R$ gives a larger convergence strip, \eqref{mld:eq:rational-kernel} continues there as a combined expression.
\end{corollary}
\begin{proof}
Partial fractions are finite, so the common strip permits termwise integration. Substitute $x=b_\nu y$ in each term. Logarithmic moments follow by differentiating in $a$, including the elementary factors $b_\nu^{a-k}$. Analytic continuation applies to the complete finite sum on any larger strip where the original integral is holomorphic.
\end{proof}

This corollary is an exact reduction procedure: denominator multiplicity, Mellin resonance, logarithmic moments, and spectral derivatives are handled by finite operations. It does not assert a comparable reduction for products of two arbitrary polylogarithms; that extension is a separate problem.


% SOURCE: sections/03-jets.tex
\section{Gamma, polygamma, and Stieltjes identities}\label{mld:sec:jets}

The same transform has two useful spectral centers. At $s=0$, the first derivative of Hurwitz zeta gives $\log\Gamma$. At $s=1$, the pole-free Hurwitz difference gives the generalized Stieltjes constants. Integer Mellin parameters then supply finite expressions in ordinary zeta jets. The resulting identities are convergent integral formulas, not finite parts.

\subsection{A log-gamma kernel and every polygamma derivative}

For $j\ge0$ and $0<\Re q<2$, set
\begin{equation}\label{mld:eq:gamma-kernel}
 \mathcal G_j(q)=\int_0^\infty
       \frac{x^{-q}\log^j x}{1+x}\Li_0^{[1]}(-x)\dd x.
\end{equation}

\begin{theorem}[Log-gamma and polygamma integral identities]\label{mld:thm:gamma-kernel}
On $0<\Re q<2$,
\begin{equation}\label{mld:eq:loggamma-integral}
 \boxed{\displaystyle
 \log\Gamma(q)=-\frac{\sin\pi q}{\pi}\mathcal G_0(q).}
\end{equation}
For every integer $r\ge1$,
\begin{equation}\label{mld:eq:polygamma-integral}
 \psi^{(r-1)}(q)=
 -\sum_{j=0}^{r}\binom rj(-1)^j\pi^{r-j-1}
       \sin\!\left(\pi q+\frac{(r-j)\pi}{2}\right)\mathcal G_j(q).
\end{equation}
All integrals are ordinary convergent integrals. In particular, $q=1$ is included without assigning a singular value to either side of \eqref{mld:eq:loggamma-integral}.
\end{theorem}

\begin{proof}
Set $a=1-q$ and $z=1$ in \eqref{mld:eq:J1}, and differentiate spectrally at $s=0$. The classical identity
\[
 \zeta^{[1]}(0,q)=\log\Gamma(q)-\tfrac12\log(2\pi)
\]
gives
\[
 \mathcal G_0(q)=-\pi\csc(\pi q)\log\Gamma(q)
\]
away from $q=1$. Multiplying by the sine and continuing gives \eqref{mld:eq:loggamma-integral} everywhere in the strip. Differentiation in $q$ gives
$\partial_q^j\mathcal G_0(q)=(-1)^j\mathcal G_j(q)$.
Apply the product rule and the formula for successive derivatives of $\sin\pi q$ to obtain \eqref{mld:eq:polygamma-integral}. Theorem~\ref{mld:thm:Mellin-master} supplies the uniform majorants for every derivative.
\end{proof}

For example,
\[
 \psi(q)=-\cos(\pi q)\mathcal G_0(q)
                  +\frac{\sin\pi q}{\pi}\mathcal G_1(q).
\]
At the half point, \eqref{mld:eq:loggamma-integral} becomes
\begin{equation}\label{mld:eq:gamma-half-integral}
 \int_0^\infty\frac{\Li_0^{[1]}(-t^2)}{1+t^2}\dd t
       =-\frac\pi4\log\pi.
\end{equation}
Here $x=t^2$ converts $\mathcal G_0(1/2)$ to twice the displayed integral.
For any two parameters $a,b$ in the strip, subtraction also gives the gamma-ratio formula
\begin{align}\label{mld:eq:gamma-ratio-integral}
 \log\frac{\Gamma(b)}{\Gamma(a)}
   =-\frac1\pi\int_0^\infty
       \frac{\sin(\pi b)x^{-b}-\sin(\pi a)x^{-a}}{1+x}
                         \Li_0^{[1]}(-x)\dd x.
\end{align}
The left side means the difference of the analytic gamma logarithms fixed in Section~\ref{mld:sec:scope}.

\subsection{Generalized Stieltjes constants without a divergent integral}

\begin{theorem}[Stieltjes difference transform]\label{mld:thm:stieltjes-kernel}
For $m\ge0$ and $0<\Re q<2$,
\begin{equation}\label{mld:eq:stieltjes-kernel}
 \boxed{\displaystyle
 \gamma_m(q)-\gamma_m
 =\frac{(-1)^{m+1}\sin\pi q}{\pi}
       \int_0^\infty\frac{x^{-q}}{1+x}\Li_1^{[m]}(-x)\dd x.}
\end{equation}
For every integer $m\ge1$,
\begin{equation}\label{mld:eq:stieltjes-zero-order}
 \int_0^\infty\frac{\Li_0^{[m]}(-x)}{x(1+x)}\dd x
       =(-1)^m m\gamma_{m-1},
\end{equation}
and equivalently
\begin{equation}\label{mld:eq:stieltjes-squared}
 \int_0^\infty\frac{\Li_1^{[m]}(-x)}{(1+x)^2}\dd x
       =(-1)^m m\gamma_{m-1}.
\end{equation}
The zeroth-derivative values of both integrals in
\eqref{mld:eq:stieltjes-zero-order}--\eqref{mld:eq:stieltjes-squared} are $-1$.
\end{theorem}

\begin{proof}
At $s=1$, the difference $\zeta(s,q)-\zeta(s)$ is entire and its $m$th derivative is $(-1)^m(\gamma_m(q)-\gamma_m)$. Apply this observation to \eqref{mld:eq:J1} at $a=1-q$. This proves \eqref{mld:eq:stieltjes-kernel}.

The integer resonance $J_1(0,s;1)=-s\zeta(s+1)$ and the convergent expansion
\begin{equation}\label{mld:eq:F-stieltjes}
 s\zeta(s+1)=1+
       \sum_{n\ge0}\frac{(-1)^n\gamma_n}{n!}s^{n+1}
\end{equation}
give \eqref{mld:eq:stieltjes-zero-order}. The case $N=2,m=1$ of
\eqref{mld:eq:integer-resonance} is $J_2(1,s;1)=-(s-1)\zeta(s)$.
Differentiate at $s=1$ to obtain \eqref{mld:eq:stieltjes-squared}. Alternatively, integration by parts directly shifts the polylogarithm order in the two integrals. All order derivatives commute with these ordinary integrals by Theorem~\ref{mld:thm:Mellin-master}.
\end{proof}

At $q=1/2$, differentiating the elementary identity
$\zeta(s,1/2)=(2^s-1)\zeta(s)$ gives
\[
 \gamma_1(1/2)-\gamma_1=-2\gamma\log2-\log^22.
\]
Thus a concrete first-order specialization of
\eqref{mld:eq:stieltjes-kernel} is
\begin{equation}\label{mld:eq:stieltjes-half-integral}
 \int_0^\infty\frac{\Li_1^{[1]}(-t^2)}{1+t^2}\dd t
       =-\frac\pi2\big(2\gamma\log2+\log^22\big).
\end{equation}
Equation \eqref{mld:eq:stieltjes-squared} supplies, for instance,
\[
 \int_0^\infty\frac{\Li_1^{[1]}(-x)}{(1+x)^2}\dd x=-\gamma,
 \qquad
 \int_0^\infty\frac{\Li_1^{[2]}(-x)}{(1+x)^2}\dd x=2\gamma_1.
\]

\subsection{The exact harmonic coefficients at spectral resonances}

The products in \eqref{mld:eq:unit-moments} are entire although their displayed zeta factor can have a pole. This is the point where finite generalized harmonic coefficients enter the Mellin calculus.

\begin{proposition}[All derivatives of a cancelled zeta pole]\label{mld:prop:spectral-cancellation}
For an integer $r\ge1$, put $s=1-r+u$. Then, as an analytic identity near $u=0$,
\begin{equation}\label{mld:eq:cancelled-zeta-germ}
 (s)_r\zeta(s+r)=(-1)^{r-1}(r-1)!
       \prod_{j=1}^{r-1}\left(1-\frac uj\right)
       \left(1+\sum_{n\ge0}\frac{(-1)^n\gamma_n}{n!}u^{n+1}\right).
\end{equation}
Consequently its $m$th spectral derivative at $s=1-r$ is
\begin{align}\label{mld:eq:cancelled-zeta-derivative}
 &(-1)^{r-1}(r-1)!m!\,[u^m]
       \left\{\prod_{j=1}^{r-1}\left(1-\frac uj\right)
       \left(1+\sum_{n\ge0}\frac{(-1)^n\gamma_n}{n!}u^{n+1}\right)\right\}.
\end{align}
Only finitely many Stieltjes constants and harmonic coefficients enter a specified derivative.
\end{proposition}
\begin{proof}
Factor the one vanishing factor from the rising factorial:
\[
 (1-r+u)_r=(-1)^{r-1}(r-1)!u
                         \prod_{j=1}^{r-1}(1-u/j).
\]
Multiply by the Laurent series for $\zeta(1+u)$. This proves the germ and its derivative formula.
\end{proof}

The polynomial coefficients are elementary symmetric functions in
$1,1/2,\ldots,1/(r-1)$. In terms of generalized harmonic numbers,
\[
 \prod_{j=1}^{r-1}(1-u/j)
 =\exp\!\left(-\sum_{k\ge1}\frac{H_{r-1}^{(k)}}{k}u^k\right).
\]
The right side is used as a formal expansion to any required finite order. Its first coefficients are
\[
 1-H_{r-1}u+
 \frac{H_{r-1}^2-H_{r-1}^{(2)}}2u^2
 -\frac{H_{r-1}^3-3H_{r-1}H_{r-1}^{(2)}+2H_{r-1}^{(3)}}6u^3+\cdots.
\]
This is a finite coefficient mechanism, distinct from interpreting an undefined expression $0\cdot\infty$ by numerical cancellation.

For example, since
$\mathcal M_1(s;1)=-\tfrac12(s)_2\zeta(s+2)$,
expansion at $s=-1$ gives
\begin{align}
 \int_0^\infty\frac{\log x\,\Li_{-1}(-x)}{x(1+x)}\dd x
      &=\frac12,\label{mld:eq:resonant-example0}\\
 \int_0^\infty\frac{\log x\,\Li_{-1}^{[1]}(-x)}{x(1+x)}\dd x
      &=\frac{\gamma-1}{2},\label{mld:eq:resonant-example1}\\
 \int_0^\infty\frac{\log x\,\Li_{-1}^{[2]}(-x)}{x(1+x)}\dd x
      &=-\gamma-\gamma_1.\label{mld:eq:resonant-example2}
\end{align}
The first identity can also be checked using the rational function
$\Li_{-1}(-x)=-x/(1+x)^2$ and a beta derivative. The spectral derivatives in the next two are different functions and should not be replaced by derivatives of that rational expression in $x$.

Higher denominator orders combine the same ingredients with the finite polynomial table. For example,
\begin{equation}\label{mld:eq:third-pole-example}
 \int_0^\infty\frac{\Li_1^{[1]}(-x)}{(1+x)^3}\dd x
       =\frac14\log(2\pi)-\frac14-\frac\gamma2.
\end{equation}
This follows by differentiating
$J_3(1,s;1)=\tfrac12(F_{s-2}(1)-F_{s-1}(1))$
at $s=1$, using $\zeta(0)=-1/2$ and
$\zeta^{[1]}(0)=-\tfrac12\log(2\pi)$.

\subsection{Antiderivatives and the existing Hurwitz-jet ladder}

The Mellin transform also produces a natural anchored antiderivative in its exponent parameter. This connects the new integral formulas to the antiderivative framework already present in the manuscript; the Hurwitz primitive itself is classical.

\begin{proposition}[Anchored primitive]\label{mld:prop:primitive}
For $-1<\Re a<1$ define
\begin{equation}\label{mld:eq:primitive-one}
 \mathcal P_s(a)=
 \frac{\zeta(s-1,1-a)-\zeta(s-1)}{s-1}-a\zeta(s),
\end{equation}
where the complete expression is continued at exceptional $s$. It is entire in $s$, vanishes at $a=0$, and satisfies
\begin{equation}\label{mld:eq:primitive-derivative}
 \partial_a\mathcal P_s(a)
   =\zeta(s,1-a)-\zeta(s)
   =-\frac{\sin\pi a}{\pi}J_1(a,s;1).
\end{equation}
At the Stieltjes center,
\begin{equation}\label{mld:eq:primitive-stieltjes}
 \left.\partial_s^m\mathcal P_s(a)\right|_{s=1}
   =\frac{\zeta^{[m+1]}(0,1-a)-\zeta^{[m+1]}(0)}{m+1}
         -a(-1)^m\gamma_m.
\end{equation}
In particular,
\begin{equation}\label{mld:eq:primitive-gamma}
 \mathcal P_1(a)=\log\Gamma(1-a)-\gamma a.
\end{equation}
\end{proposition}
\begin{proof}
The Hurwitz parameter derivative
$\partial_q\zeta(v,q)=-v\zeta(v+1,q)$ proves
\eqref{mld:eq:primitive-derivative} away from the displayed singularities. The right side is entire in $s$, and its integral from $0$ to $a$ gives an entire continuation of \eqref{mld:eq:primitive-one}. At $s=1+u$, the numerator has constant term
$\zeta(0,1-a)-\zeta(0)=a$. Its quotient contributes $a/u$, cancelling the pole of $a\zeta(1+u)$. Comparing the remaining coefficients gives \eqref{mld:eq:primitive-stieltjes}, and the log-gamma identity gives \eqref{mld:eq:primitive-gamma}.
\end{proof}

For completeness, every positive integer primitive order $r$ has a finite anchored expression:
\begin{align}\label{mld:eq:primitive-r}
 \mathcal P_{s,r}(a)
  ={}&\frac{\displaystyle\zeta(s-r,1-a)
       -\sum_{j=0}^{r-1}\frac{(s-r)_j\zeta(s-r+j)}{j!}a^j}
                   {(s-r)_r}
        -\frac{a^r}{r!}\zeta(s).
\end{align}
Away from the exceptional spectral parameters, $r$ differentiations in $a$ give \eqref{mld:eq:primitive-derivative}; all lower initial derivatives vanish. Equivalently,
\[
 \mathcal P_{s,r}(a)=\frac1{(r-1)!}\int_0^a
       (a-v)^{r-1}\big[\zeta(s,1-v)-\zeta(s)\big]\dd v.
\]
This integral proves that all apparent spectral singularities in the combined expression \eqref{mld:eq:primitive-r} are removable. Taylor extraction at those singularities produces the complete homogeneous harmonic coefficients familiar from the manuscript's negative-polygamma ladder. The anchor fixes the integration polynomial, so no unrecorded integration constants enter.

Finally, argument differentiation and logarithmic integration commute with the Mellin operation:
\begin{equation}\label{mld:eq:argument-ladder}
 z\partial_z J_N(a,s;z)=J_N(a,s-1;z),\qquad
 \int \frac{J_N(a,s;z)}{z}\dd z=J_N(a,s+1;z)+C.
\end{equation}
The primitive is local on a simply connected part of $\mathcal S$. This follows directly from the polylogarithm order ladder and the compact bounds proved earlier.


% SOURCE: sections/04-harmonic.tex
\section{A two-shift reflection for generalized harmonic sums}
\label{mld:gh:sec:reflection}\label{mld:sec:harmonic}

The generalized harmonic order and the elementary-symmetric harmonic index
in the repository are different parameters. We use new notation to keep
them separate. Put
\[
 h_n^{(r)}(u)=\sum_{j=1}^n\frac1{(j-u)^r},\qquad
 Q_{p,r}=\sum_{n=1}^{\infty}\frac{(-1)^nH_n^{(r)}}{(2n+1)^p},
 \qquad p,r\geq1.
\]
Thus $h_n^{(r)}(0)=H_n^{(r)}$, and $Q_{p,1}=S_p$ in the manuscript.
In contrast, the manuscript's $S_{p,r}$ uses
$e_r(n)=[v^r]\prod_{j=1}^n(1+v/j)$; it is not $Q_{p,r}$.

The unshifted even-weight evaluations belong to the established theory
of linear Euler $R$-sums. In the normalization of Xu--Wang
\cite[Section~3.4, equation~(20)]{XuWang},
\begin{equation}\label{mld:gh:eq:normalization}
 R_{r,p}^{1,-1}=2^p Q_{p,r}.
\end{equation}
We derive a two-variable digamma resummation, its arbitrary mixed
derivatives, and a concrete rational-shift evaluation. The scalar parity
theorem is credited to that literature. The priority of the particular
two-variable formulation below has not been established.

\subsection{Ordinary convergence, including the first outer power}

Define
\begin{align}
 A(z,u)&=\sum_{n=1}^{\infty}
       \frac{(-1)^n h_n^{(1)}(u)}{2n+1-z},\label{mld:gh:eq:A}\\
 D(v)&=\frac14\left\{
       \psi\left(\frac{v+3}{4}\right)
       -\psi\left(\frac{v+1}{4}\right)\right\}.
       \label{mld:gh:eq:D}
\end{align}
Initially, $D(v)=\int_0^1x^v/(1+x^2)\,dx$ for $\Re v>-1$;
the displayed expression supplies its meromorphic continuation.
The first sum in \eqref{mld:gh:eq:A} is an ordinary convergent alternating
series.

\begin{lemma}[Local uniform convergence]\label{mld:gh:lem:convergence}
For positive integers $p,r$, the series
\[
 \sum_{n=1}^{\infty}
 \frac{(-1)^n h_n^{(r)}(u)}{(2n+1-z)^p}
\]
converges locally uniformly on
\[
 \Omega_+=
 \bigl(\mathbb C\setminus\{3,5,7,\ldots\}\bigr)
 \times\bigl(\mathbb C\setminus\{1,2,3,\ldots\}\bigr).
\]
It is holomorphic there and can be differentiated term by term any
finite number of times in $z$ and $u$. It is absolutely convergent
when $p\geq2$.
\end{lemma}
\begin{proof}
Fix a compact subset of $\Omega_+$. Uniformly on that subset,
$h_n^{(1)}(u)=O(1+\log n)$, while $h_n^{(r)}(u)=O(1)$ for $r\geq2$.
For
$f_n=h_n^{(r)}(u)(2n+1-z)^{-p}$, the exact difference is
\begin{align*}
 f_{n+1}-f_n
 ={}&\frac{1}{(n+1-u)^r(2n+3-z)^p}\\
 &+h_n^{(r)}(u)
       \left\{(2n+3-z)^{-p}-(2n+1-z)^{-p}\right\}.
\end{align*}
Consequently $f_n\to0$ uniformly and
$f_{n+1}-f_n=O(n^{-p-r}+(1+\log n)n^{-p-1})$ uniformly.
The sum of these differences is absolutely convergent, so grouping
successive even and odd indices proves local uniform convergence of
the original alternating series. Its individual terms are holomorphic.
The same estimates apply to every termwise derivative, because
derivatives increase $p$ or $r$. This proves the assertion, including
the conditionally convergent case $p=1$. The stated absolute convergence
follows directly from the bounds for $h_n^{(r)}$.
\end{proof}

\subsection{The exact bivariate identity}

\begin{theorem}[Two-shift reflection]\label{mld:gh:thm:reflection}
On
\[
 \Omega=
 \bigl(\mathbb C\setminus\{\pm3,\pm5,\pm7,\ldots\}\bigr)
 \times\bigl(\mathbb C\setminus(\mathbb Z\setminus\{0\})\bigr),
\]
one has
\begin{equation}\label{mld:gh:eq:reflection}
 \boxed{A(z,u)+A(-z,-u)=B(z,u),}
\end{equation}
where
\begin{align}
 B(z,u)={}&-\frac{\pi}{2\cos(\pi z/2)}
   \left\{\psi(1-u)-\psi\left(\frac{1+z}{2}-u\right)\right\}
       \notag\\
 &+\frac{\pi D(z-2u)}{\sin\pi u}-\frac{D(z)}{u}.
       \label{mld:gh:eq:B}
\end{align}
Apparent singularities of the complete expression on $\Omega$ are
removable. This convention applies in particular at $u=0$, at $z=\pm1$,
and where an individual digamma term in \eqref{mld:gh:eq:B} has a pole.
There is no assertion that its summands can be evaluated separately
at such points.

For $|z|<3$ and $|u|<1$, the ordinary Taylor expansion of $A$ is
\begin{equation}\label{mld:gh:eq:Agenerating}
 A(z,u)=\sum_{p,r\geq1} Q_{p,r}\,z^{p-1}u^{r-1}.
\end{equation}
Thus $B$ determines precisely the terms with $p+r$ even in this germ.
\end{theorem}

\begin{proof}
We first prove the identity near $(z,u)=(0,0)$. Set
\[
 G_r(z)=\int_0^1
      \frac{x^{-z}\operatorname{Li}_r(-x^2)}{1+x^2}\,dx.
\]
The harmonic generating function
$\operatorname{Li}_r(w)/(1-w)=\sum_{n\geq1}H_n^{(r)}w^n$
gives
$G_r(z)=\sum_{p\geq1}Q_{p,r}z^{p-1}$.
One can first insert a common Abel factor and then pass to its limit;
the preceding lemma identifies that limit with the ordinary sum.
Alternatively, the logarithmic moment representation is
\begin{equation}\label{mld:gh:eq:moment}
 Q_{p,r}=\frac1{(p-1)!}\int_0^1
   \frac{(-\log x)^{p-1}\operatorname{Li}_r(-x^2)}{1+x^2}\,dx.
\end{equation}
For $0\leq y\leq1$ and integer $r\geq1$,
$|\operatorname{Li}_r(-y)|\leq y$ by the alternating-series test.
It follows from \eqref{mld:gh:eq:moment} that
$|Q_{p,r}|\leq3^{-p}$. Hence the double power series in
\eqref{mld:gh:eq:Agenerating} converges normally for the stated radii.
Its coefficients also follow by the justified derivatives of
\eqref{mld:gh:eq:A} at the origin.

The full Mellin identity, after the substitution $y=x^2$, is
\begin{equation}\label{mld:gh:eq:fullMellin}
 M_r(z):=\int_0^\infty
     \frac{x^{-z}\operatorname{Li}_r(-x^2)}{1+x^2}\,dx
 =-\frac{\pi}{2\cos(\pi z/2)}
   \left\{\zeta\left(r,\frac{1+z}{2}\right)-\zeta(r)\right\}.
\end{equation}
It initially suffices to take $|\Re z|<1$. At $r=1$, the expression
in braces denotes the finite difference
$\psi(1)-\psi((1+z)/2)$.

For $L\in\mathbb R$, define the polynomial
\begin{equation}\label{mld:gh:eq:Pr}
 P_r(L)=\frac{(2L)^r}{r!}
     +2\sum_{j=1}^{\lfloor r/2\rfloor}
       \eta(2j)\frac{(2L)^{r-2j}}{(r-2j)!},
 \qquad \eta(s)=(1-2^{1-s})\zeta(s).
\end{equation}
Then
\begin{equation}\label{mld:gh:eq:inversion}
 \operatorname{Li}_r(-e^{2L})
   +(-1)^r\operatorname{Li}_r(-e^{-2L})=-P_r(L).
\end{equation}
Here the signs and branch convention can be checked without invoking
a complex inversion formula. The derivative of the left side is twice
the corresponding expression for $r-1$. At $r=0$ that expression
equals $-1$. At $L=0$ it is zero for odd $r$, and $-2\eta(r)$ for
even positive $r$. These data give exactly \eqref{mld:gh:eq:Pr}.

Splitting \eqref{mld:gh:eq:fullMellin} at $1$ and substituting $x\mapsto1/x$
in its second part gives
\begin{equation}\label{mld:gh:eq:integerreflection}
 G_r(z)-(-1)^rG_r(-z)
 =M_r(z)+(-1)^r\int_0^1\frac{x^zP_r(\log x)}{1+x^2}\,dx.
\end{equation}
Now sum with multiplier $u^{r-1}$ over $r\geq1$. In a sufficiently
small neighborhood, for example $|z|+2|u|<1$, all the integrations and
sums involved are dominated by integrable powers of $x$ at their
endpoints. The two elementary generating identities are
\begin{align}
 \sum_{r\geq1}u^{r-1}\{\zeta(r,b)-\zeta(r)\}
   &=\psi(1-u)-\psi(b-u),\label{mld:gh:eq:zetasum}\\
 \sum_{r\geq0}P_r(L)u^r
   &=\frac{\pi u}{\sin\pi u}e^{2uL},
       \qquad P_0=1.\label{mld:gh:eq:Psum}
\end{align}
In \eqref{mld:gh:eq:zetasum}, the $r=1$ value is interpreted as above;
the formula follows by summing
$(n+b-u)^{-1}-(n+1-u)^{-1}$ over $n\geq0$.
Equation \eqref{mld:gh:eq:Psum} follows from
$\pi u/\sin\pi u=1+2\sum_{j\geq1}\eta(2j)u^{2j}$.

The sum of the left sides of \eqref{mld:gh:eq:integerreflection} is
$A(z,u)+A(-z,-u)$. Equation \eqref{mld:gh:eq:zetasum} gives the first
line of \eqref{mld:gh:eq:B}; equation \eqref{mld:gh:eq:Psum}, with $u$
replaced by $-u$, gives
\[
 \int_0^1\frac{x^z}{1+x^2}
   \left\{\frac{\pi x^{-2u}}{\sin\pi u}-\frac1u\right\}\,dx
 =\frac{\pi D(z-2u)}{\sin\pi u}-\frac{D(z)}u.
\]
This proves the claimed identity as an analytic germ.

The paired series in Lemma~\ref{mld:gh:lem:convergence} gives a meromorphic
continuation of $A$ to $\mathbb C^2$: near any candidate polar
hyperplane, multiply by its local linear factors; the same uniform
tail estimates apply. Its possible polar hyperplanes are
$z=3,5,\ldots$ and $u=1,2,\ldots$.
The right side of \eqref{mld:gh:eq:B} is also meromorphic on $\mathbb C^2$.
The identity theorem for meromorphic functions extends the equality
from the germ to the whole connected domain. The left side is
holomorphic on $\Omega$, so all apparent poles of the complete
right side there cancel. Finally, the sign in the Taylor projection
is $(-1)^{(p-1)+(r-1)}=(-1)^{p+r}$, proving the final assertion.
\end{proof}

\begin{corollary}[Every reflected shifted harmonic sum]
\label{mld:gh:cor:derivatives}
For positive integers $p,r$ and $(z,u)\in\Omega$,
\begin{align}
 &\sum_{n=1}^{\infty}(-1)^n
   \left\{\frac{h_n^{(r)}(u)}{(2n+1-z)^p}
       +(-1)^{p+r}\frac{h_n^{(r)}(-u)}{(2n+1+z)^p}\right\}
       \notag\\
 &\hspace{14mm}=
 \frac{\partial_z^{p-1}\partial_u^{r-1}B(z,u)}
      {(p-1)!(r-1)!}.\label{mld:gh:eq:allmixed}
\end{align}
The right side is a finite expression in polygamma functions and
elementary trigonometric functions. The equality holds for ordinary
conditionally convergent sums when $p=1$.
\end{corollary}
\begin{proof}
Differentiate \eqref{mld:gh:eq:reflection}. Lemma~\ref{mld:gh:lem:convergence}
justifies the operation, including $p=1$. The identities
$\partial_u^{r-1}h_n^{(1)}(u)=(r-1)!h_n^{(r)}(u)$ and
$\partial_z^{p-1}(2n+1-z)^{-1}=(p-1)!(2n+1-z)^{-p}$
give the first term. Applying the chain rule to $A(-z,-u)$ contributes
the factor $(-1)^{p+r}$.
\end{proof}

At rational parameter pairs in the domain, the right side uses only
polygamma values at rational arguments and elementary constants. The
digamma differences reduce by Gauss's formula to logarithms of algebraic
numbers and algebraic multiples of $\pi$. Positive polygamma orders
reduce to Hurwitz zeta values, and hence by residue-class filters to
depth-one polylogarithms at roots of unity. This gives an explicit
depth-one expression for the reflected combination; it does not assert
the same for either summand separately.

\subsection{A finite formula for the unshifted parity coefficients}

Write $\beta(s)=\sum_{n\ge0}(-1)^n/(2n+1)^s$ for the Dirichlet
beta function, initially on $\Re s>0$ and continued analytically.
For $j\geq0$, define
\[
 d_0(r)=
 \begin{cases}2\log2,&r=1,\\(2^r-2)\zeta(r),&r>1,\end{cases}
 \qquad
 d_j(r)=\frac{(r)_j(2^{r+j}-1)\zeta(r+j)}{2^j j!}
 \quad(j\geq1),
\]
and let
$c_{2\ell}=|E_{2\ell}|\pi^{2\ell}/(2^{2\ell}(2\ell)!)$, where
$\sec t=\sum_{\ell\geq0}|E_{2\ell}|t^{2\ell}/(2\ell)!$.
Put
\begin{align}
 C_{p,r}={}&2^r\binom{p+r-1}{r}\beta(p+r)\notag\\
 &+2\sum_{j=1}^{\lfloor r/2\rfloor}
      2^{r-2j}\eta(2j)
      \binom{p+r-2j-1}{r-2j}\beta(p+r-2j).
      \label{mld:gh:eq:Cpr}
\end{align}
The coefficient form of the reflection is
\begin{equation}\label{mld:gh:eq:finiteparity}
 \boxed{
 Q_{p,r}=\frac{(-1)^{r+1}}2
 \left\{C_{p,r}-\frac\pi2
   \sum_{\ell=0}^{\lfloor(p-1)/2\rfloor}
        c_{2\ell}d_{p-1-2\ell}(r)\right\},
 \qquad p+r\ \text{even}.}
\end{equation}
Indeed, extract $z^{p-1}$ in \eqref{mld:gh:eq:integerreflection}; the
coefficient on its left is $2Q_{p,r}$ in the indicated parity.
Taylor expansion of the Hurwitz zeta gives the $d_j$ terms.
Integration of each monomial in \eqref{mld:gh:eq:Pr} gives the beta
terms in \eqref{mld:gh:eq:Cpr}; the common sign is $(-1)^{r+1}$.
This is a finite coefficient algorithm for every $p,r$, with no
integer-relation search.

For example, writing $G=\beta(2)$,
\begin{align}
 Q_{2,2}&=\frac{7\pi}{4}\zeta(3)-6\beta(4)-\frac{\pi^2G}{12},
       \label{mld:gh:eq:Q22}\\
 Q_{1,3}&=4\beta(4)+\frac{\pi^2G}{6}
                     -\frac{3\pi}{2}\zeta(3),\\
 Q_{4,2}&=\frac{31\pi}{8}\zeta(5)+\frac{7\pi^3}{32}\zeta(3)
                     -20\beta(6)-\frac{\pi^2}{12}\beta(4),\\
 Q_{3,3}&=40\beta(6)+\frac{\pi^2}{2}\beta(4)
                     -\frac{93\pi}{8}\zeta(5)
                     -\frac{3\pi^3}{16}\zeta(3),\\
 Q_{4,4}&=-280\beta(8)-\frac{10\pi^2}{3}\beta(6)
                     -\frac{7\pi^4}{720}\beta(4)
                     +\frac{635\pi}{8}\zeta(7)
                     +\frac{31\pi^3}{16}\zeta(5).
\end{align}
Every term has total weight $p+r$. Only odd zeta values, even beta
values, $\pi$, and, in the $r=1$ case, $\log2$ occur.

\subsection{A rational-shift identity with an elementary answer}

\begin{corollary}\label{mld:gh:cor:quartershift}
The ordinary convergent nested sum satisfies
\begin{equation}\label{mld:gh:eq:quartershift}
 \boxed{
 \sum_{n=1}^{\infty}\frac{(-1)^n}{2n+1}
  \sum_{j=1}^n\left\{\frac1{j-\frac14}+\frac1{j+\frac14}\right\}
 =\pi\bigl(\log(1+\sqrt2)-1\bigr).}
\end{equation}
Equivalently, dividing by $32$, the inner summand is
$j/(16j^2-1)$ and the right side is
$\pi(\log(1+\sqrt2)-1)/32$.
\end{corollary}
\begin{proof}
Set $(z,u)=(0,1/4)$ in \eqref{mld:gh:eq:reflection}.
The digamma reflection formula gives
$\psi(3/4)-\psi(1/4)=\pi$, and $D(0)=\pi/4$.
The remaining resolvent is
\[
 D(-1/2)=2\int_0^1\frac{dt}{1+t^4}
 =\frac{\pi}{2\sqrt2}
       +\frac{\log(1+\sqrt2)}{\sqrt2}.
\]
The last integral follows by factoring
$1+t^4=(t^2+\sqrt2t+1)(t^2-\sqrt2t+1)$ and integrating the two
partial fractions. Substitution into \eqref{mld:gh:eq:B} gives
$-\pi^2/2+\pi\sqrt2D(-1/2)-\pi$, which is the stated answer.
\end{proof}

\subsection{Verification and the unresolved odd-weight direction}

The accompanying \texttt{derive\_parity.py} generates the finite symbolic
formula \eqref{mld:gh:eq:finiteparity}. It independently compares all $25$
ordered pairs with $p+r\leq10$ and $p+r$ even against real quadrature
of \eqref{mld:gh:eq:moment} at $60$ decimal working digits. The largest
absolute discrepancy is less than $3\times10^{-58}$.
Three separate $450$-term Euler evaluations of \eqref{mld:gh:eq:A}, including
a complex parameter pair and a point outside the initial small-parameter
proof neighborhood, agree with \eqref{mld:gh:eq:B} to more than $115$
decimal places. The JSON file records every tested pair and residual.
Three mixed-derivative checks, including $p=1,r=2$, and the quarter-shift
specialization are checked by the same independent series method; their
absolute discrepancies are below $3\times10^{-136}$.
These checks audit the signs and normalizations; the arbitrary-order
proof is the analytic argument above.

At $(z,u)=(0,0)$, equation \eqref{mld:gh:eq:allmixed} has coefficient
$1+(-1)^{p+r}$. It therefore gives no value for $Q_{p,r}$ when
$p+r$ is odd. In particular, $S_6=Q_{6,1}$ and $S_8=Q_{8,1}$ remain
outside the scalar parity projection. This is a precise limitation
of the reflection method, not a proof of arithmetic irreducibility.
The repository already exhibits a nonzero separating functional for
the restricted depth-two product presentation of the $S_6$ candidate;
the present identity does not supply the missing odd-weight relation.


% SOURCE: sections/05-dilation.tex
% Self-contained contribution; standard amsmath/amssymb/amsthm environments.
% Suggested bibliography keys: incoming-shifted, incoming-convolution,
% incoming-closure, dlmf-hurwitz.  All statements use ambient x cutoffs.
\section{Dilation and unequal-frequency correlations}
\label{mld:dil:section}\label{mld:sec:dilation}

Two explicit questions in the incoming reports concern the effect of
integer dilation on canonical periodic finite parts, and products with
unequal integer dilations.  They occur in the further-research sections
of the shifted-Hurwitz and convolution-calculus deliveries.  This section
answers those questions for every nonnegative Stieltjes index and every
nonnegative parameter derivative order, provided the two sets of
singular points are disjoint.  The location of the surviving ordinary
special-function values is determined by the reduced dilation ratio;
their coefficients also remember the physical coordinate used for
the finite part.

\subsection{Conventions and the elementary contact identity}
Write $\mathbb T=\mathbb R/\mathbb Z$, and use normalized Lebesgue
measure.  The spectral convention is
\[
 \zeta(1-z,x)=-\frac1z+
       \sum_{n\ge0}\gamma_n(x)\frac{z^n}{n!},
 \qquad \gamma_0(x)=-\psi(x),
 \qquad 0<x<1.
\]
For a smooth periodic test function $\phi$, let $T_{n,r}$ be the
constant-term extension of $\gamma_n^{(r)}(x)$ from $(0,1)$:
\[
 \langle T_{n,r},\phi\rangle
 =\operatorname{CT}_{\epsilon\downarrow0}
      \int_\epsilon^1\gamma_n^{(r)}(x)\phi(x)\,dx.
\]
Here $\operatorname{CT}$ retains the constant term after all negative
powers of $\epsilon$ multiplied by powers of $\log\epsilon$, and all
positive powers of $\log\epsilon$, have been removed.  No finite
point-supported term is added.  Thus, with $G_n=T_{n,0}$,
\[
 \langle G_n,\phi\rangle
 =\int_0^1\left\{\gamma_n(x)\phi(x)
                  -\frac{\log^n x}{x}\phi(0)\right\}\,dx.
\]
The shift identity $\gamma_n(x)=x^{-1}\log^n x+\gamma_n(1+x)$
proves convergence and the equality of these definitions.

Put
\begin{equation}
 p_r(z)=\prod_{j=1}^r\left(1-\frac zj\right),\qquad
 h_{r,n}=(-1)^n n!e_{n+1}(1,1/2,\ldots,1/r),
 \label{mld:dil:p-h}
\end{equation}
where the empty product is one and elementary symmetric functions of
degree greater than $r$ are zero.  The undilated contact identity, proved
in the incoming convolution and shifted-Hurwitz reports
\cite{IncomingConvolution,IncomingShifted}, is
\begin{equation}
 D^rG_n=T_{n,r}-h_{r,n}\delta_0^{(r)}.
 \label{mld:dil:unit-contact}
\end{equation}
For completeness, its origin can be read directly from the meromorphic
periodic Hurwitz family.  If
\[
 \widehat R_z(0)=0,\qquad
 \widehat R_z(k)=(2\pi i k)^{-z}\quad(k\ne0),
 \qquad Z(1-z)=\Gamma(z)R_z,
\]
then
\[
 Z(1-z)=\frac{\delta_0-1}{z}
            +\sum_{n\ge0}G_n\frac{z^n}{n!}.
\]
For $r\ge1$, the local singular part of the pointwise $r$th derivative
is $(-1)^r r!p_r(z)x^{z-r-1}$.  The meromorphic extension of
$x^{z-r-1}$ has residue $(-1)^r\delta_0^{(r)}/r!$ at $z=0$.
Consequently its polar contribution is
$p_r(z)\delta_0^{(r)}/z$.  Comparing coefficients with
$D^rZ(1-z)$ gives \eqref{mld:dil:unit-contact}, since
$n![z^{n+1}]p_r(z)=-h_{r,n}$.
The case $r=0$ is immediate.  This derivation also specifies all signs.
The coefficients can equally be written through generalized harmonic
numbers: for $|z|<1$,
\[
 p_r(z)=\exp\!\left(-\sum_{j\ge1}\frac{H_r^{(j)}}jz^j\right),
 \qquad H_r^{(j)}=\sum_{\ell=1}^r\ell^{-j}.
\]
For example $e_1=H_r$ and $e_2=(H_r^2-H_r^{(2)})/2$.

\subsection{The exact dilation anomaly}
For a positive integer $q$, let $U_q$ be pullback by the covering map
$x\mapsto qx$:
\[
 \langle U_qT,\phi\rangle
 =\left\langle T,\frac1q\sum_{j=0}^{q-1}
                         \phi\!\left(\frac{x+j}{q}\right)\right\rangle.
\]
For an ordinary function this is $U_qf(x)=f(qx)$.  Write
\[
 \Delta_q=\sum_{j=0}^{q-1}\delta_{j/q},\qquad
 U_q\delta_0^{(r)}=q^{-r-1}\Delta_q^{(r)}.
\]
Let $X_{q,n,r}$ be the constant-term extension of
$\gamma_n^{(r)}(\{qx\})$ using the \emph{ambient coordinate $x$}:
at each grid point $j/q$, delete the right interval
$(j/q,j/q+\epsilon)$ and take the constant term as $\epsilon$ tends to
zero. The transported finite part $U_qT_{n,r}$ takes the constant term
in a source cutoff $\eta$, corresponding to ambient cutoff $\eta/q$.
The ambient finite part instead takes the constant term in $\epsilon$
after the source lower limit becomes $q\epsilon$. The variable in which
the constant term is extracted is part of the normalization.

\begin{theorem}[All-order finite-part dilation law]
\label{mld:dil:scale-thm}
For $q\ge1$ and $n,r\ge0$, define the explicit polynomial
\begin{equation}
 c_{n,r}(L)=n![z^{n+1}]p_r(z)(e^{Lz}-1)
 =n!\sum_{j=0}^{\min(r,n)}
       \frac{(-1)^j e_j(1,1/2,\ldots,1/r)}{(n+1-j)!}
       L^{n+1-j}.
 \label{mld:dil:c}
\end{equation}
Then
\begin{equation}
 \boxed{X_{q,n,r}=U_qT_{n,r}
       -q^{-r-1}c_{n,r}(\log q)\Delta_q^{(r)}.}
 \label{mld:dil:scale-law}
\end{equation}
In particular $c_{n,0}(L)=L^{n+1}/(n+1)$ and $c_{0,r}(L)=L$.
\end{theorem}

\begin{proof}
Near a grid point use $t=x-j/q>0$.  The singular part of the generating
function for the argument derivatives is
\[
 (-1)^r r!p_r(z)(qt)^{z-r-1}
 =(-1)^r r!q^{-r-1}p_r(z)e^{z\log q}t^{z-r-1}.
\]
After pairing with a test function, only its Taylor term of degree $r$
produces a pole at $z=0$.  Its contribution is
\[
 q^{-r-1}\frac{p_r(z)e^{z\log q}}z\delta_{j/q}^{(r)}.
\]
The pullback of the source-coordinate family instead has polar
contribution $q^{-r-1}p_r(z)\delta_{j/q}^{(r)}/z$.
The regular Taylor coefficients left after subtracting these complete
polar contributions are respectively the ambient-coordinate constant
terms and the pulled-back constant terms.  Their difference is therefore
the negative of
\[
 q^{-r-1}p_r(z)\frac{e^{z\log q}-1}{z}\delta_{j/q}^{(r)}.
\]
Extract $z^n/n!$ and sum over the grid.  When $r=0$, the scalar spectral
pole $-1/z$ occurs in both calculations and cancels.  Subtracting the
test-function Taylor polynomial through degree $r$ makes every
remaining local integral holomorphic near $z=0$, which justifies the
coefficient comparison.  Expanding the finite product gives
\eqref{mld:dil:c}.
\end{proof}

\begin{corollary}[Differentiation and dilation with fixed ambient scale]
Define
\begin{equation}
 b_{r,n}(L)=\sum_{j=0}^n\binom nj L^{n-j}h_{r,j}.
 \label{mld:dil:b}
\end{equation}
Then
\begin{equation}
 \boxed{D^rX_{q,n,0}=q^rX_{q,n,r}
          -\frac{b_{r,n}(\log q)}q\Delta_q^{(r)}.}
 \label{mld:dil:commutator}
\end{equation}
For every $n\ge0$ this specializes to
\begin{equation}
 DX_{q,n,0}=qX_{q,n,1}
             -\frac{(\log q)^n}{q}\Delta_q',
 \label{mld:dil:first-commutator}
\end{equation}
where $(\log1)^0=1$.
\end{corollary}
\begin{proof}
Differentiate \eqref{mld:dil:scale-law} for $r=0$, use
$D^rU_q=q^rU_qD^r$ and \eqref{mld:dil:unit-contact}, and compare with
\eqref{mld:dil:scale-law} at general $r$.  The remaining scalar identity is
\[
 b_{r,n}(L)=h_{r,n}+\frac{L^{n+1}}{n+1}-c_{n,r}(L)
          =n![z^{n+1}](1-p_r(z))e^{Lz}.
\]
For $r=1$, only $h_{1,0}=1$ is nonzero.
\end{proof}

The undilated anomaly vanishes when $n\ge r$.  Equation
\eqref{mld:dil:first-commutator} shows that this vanishing generally fails
after ambient-coordinate dilation: already $n=r=1$ gives
$DX_{q,1,0}=qX_{q,1,1}-(\log q)\Delta_q'/q$.
This is an extension of the incoming contact theorem, not a correction
to its correctly normalized undilated statement.

The scale corrections are compatible with successive dilations.  The
following polynomial identities hold for all $L,M\in\mathbb C$:
\begin{align}
 c_{n,r}(L+M)&=c_{n,r}(M)
       +\sum_{j=0}^n\binom nj M^{n-j}c_{j,r}(L),\notag\\
 b_{r,n}(L+M)&=\sum_{j=0}^n\binom nj M^{n-j}b_{r,j}(L).
 \label{mld:dil:scale-composition}
\end{align}
Indeed, use $e^{(L+M)z}-1=(e^{Mz}-1)+e^{Mz}(e^{Lz}-1)$ in
\eqref{mld:dil:c}, and use the last generating expression in the preceding
proof for $b$.  Thus a two-step change of local scale gives the same
finite correction as the combined scale; no independent finite
constant is introduced by composing the operations.

\subsection{A multiplication theorem including the point masses}
Let $\tau_aT(x)=T(x+a)$ and $L=\log q$.
\begin{theorem}[Distributional multiplication law]
For $n\ge0$,
\begin{equation}
 \sum_{j=0}^{q-1}\tau_{j/q}G_n
 =q\sum_{i=0}^n\binom ni(-L)^{n-i}U_qG_i
     +\frac{(-L)^{n+1}}{n+1}(\Delta_q-q).
 \label{mld:dil:multiplication}
\end{equation}
More generally, for $r\ge1$,
\begin{equation}
 \sum_{j=0}^{q-1}\tau_{j/q}T_{n,r}
 =q^{r+1}\sum_{i=0}^n\binom ni(-L)^{n-i}U_qT_{i,r}
                  +c_{n,r}(-L)\Delta_q^{(r)}.
 \label{mld:dil:derivative-multiplication}
\end{equation}
\end{theorem}
\begin{proof}
The classical Hurwitz multiplication formula gives, as meromorphic
periodic distributions,
\[
 \sum_{j=0}^{q-1}\tau_{j/q}Z(1-z)
                  =q^{1-z}U_qZ(1-z).
\]
It can also be proved directly from the Fourier coefficients: the
sum annihilates frequencies not divisible by $q$.  Expanding the
distributional Laurent series proves \eqref{mld:dil:multiplication}; the
residue transforms as $\sum\tau_{j/q}(\delta_0-1)=\Delta_q-q$.

Differentiate the pointwise multiplication formula $r$ times and then
take ambient-coordinate constant terms.  Substitution of
\eqref{mld:dil:scale-law} reduces the additional coefficient to
\[
 -\sum_{i=0}^n\binom ni(-L)^{n-i}c_{i,r}(L)
 =n![z^{n+1}]p_r(z)(e^{-Lz}-1)=c_{n,r}(-L).
\]
This gives \eqref{mld:dil:derivative-multiplication}.
\end{proof}

\subsection{The undilated kernel and an explicit finite coefficient rule}
\label{mld:dil:undilated-algorithm}

The unequal-dilation theorem below uses a known undilated closure from
the incoming reports \cite{IncomingClosure,IncomingShifted}. We record its
generating function and derivation so that the resulting algorithm is
self-contained. These undilated statements are prior results within the
supplied material.

For $0<t<1$ and $\Re s,\Re v<1$, set
\[
 C(s,v;t)=\int_0^1\zeta(s,x)\zeta(v,\{x+t\})\dd x.
\]
The integral is ordinary: the two singularities are separated. Its two
equivalent continuations are
\begin{align}
 C(s,v;t)={}&\Gamma(1-s)\Gamma(1-v)(2\pi)^{s+v-2}
 \bigl\{e^{-\pi\ii(s-v)/2}\Li_{2-s-v}(e^{2\pi\ii t})\notag\\
 &\hspace{38mm}+e^{\pi\ii(s-v)/2}\Li_{2-s-v}(e^{-2\pi\ii t})\bigr\},
 \label{mld:dil:C-polylog}\\
 C(s,v;t)={}&\Gamma(s+v-1)
 \left\{\frac{\Gamma(1-s)}{\Gamma(v)}\zeta(s+v-1,t)
       +\frac{\Gamma(1-v)}{\Gamma(s)}\zeta(s+v-1,1-t)\right\}.
 \label{mld:dil:C-hurwitz}
\end{align}
All apparent singularities are interpreted in the complete expression.
To prove the first formula, start with $\Re s,\Re v<0$ and the absolutely
convergent Hurwitz Fourier series, whose nonzero coefficients are
\[
 \widehat\zeta_s(k)=\Gamma(1-s)(2\pi\ii k)^{s-1},\qquad k\ne0.
\]
The zero coefficient is zero. Integrating the product retains the
opposite frequencies and yields \eqref{mld:dil:C-polylog}. For the second
formula, insert the same Fourier expansion with spectral order $s+v-1$
on its right. Euler's gamma reflection identity reduces the coefficient
of the positive-argument polylogarithm by
\[
 \frac{\sin(\pi s)e^{-\pi\ii(s+v)/2}
       +\sin(\pi v)e^{\pi\ii(s+v)/2}}{\sin\pi(s+v)}
       =e^{-\pi\ii(s-v)/2}.
\]
The other coefficient is obtained by reversing the signs in the phases.
This calculation holds off the integer hyperplanes; analytic continuation
extends it across their removable singularities. The local decomposition
$\zeta(s,x)=x^{-s}+\zeta(s,1+x)$ justifies joint holomorphy of the
integral on $\Re s,\Re v<1$, and hence the continuation to that region.
This is the classical Fourier and product-integral mechanism
\cite{DLMFhurwitz,EspinosaMoll} in the circular-translation convention.

Define, near $(u,v)=(0,0)$,
\begin{align*}
 \mathcal A(u,v)&=
  \frac{\Gamma(1+u+v)\Gamma(1-v)}{\Gamma(1+u)},
 &\mathcal B(u,v)&=\mathcal A(v,u),\\
 \mathcal K(u,v;t)&=uv\,C(1+u,1+v;t).
\end{align*}
Equation~\eqref{mld:dil:C-hurwitz} gives
\begin{equation}
 \mathcal K(u,v;t)=
 -u\mathcal A(u,v)\zeta(1+u+v,1-t)
 -v\mathcal B(u,v)\zeta(1+u+v,t).
 \label{mld:dil:K-hurwitz}
\end{equation}
This function is holomorphic at the origin. Indeed,
$\mathcal A(u,-u)=\mathcal B(u,-u)=1$, so
$(u\mathcal A+v\mathcal B)/(u+v)$ has a removable denominator. Removing
the common Hurwitz pole in \eqref{mld:dil:K-hurwitz} leaves the negative of that
quotient plus analytic Stieltjes series.

For the canonical undilated finite part one consequently has the
terminating coefficient rule
\begin{equation}
 \boxed{I_{mn}(t)=(-1)^{m+n}m!n!
       [u^{m+1}v^{n+1}]\mathcal K(u,v;t).}
 \label{mld:dil:I-coefficient}
\end{equation}
Here $I_{mn}$ denotes the constant term of
$\int\gamma_m(x)\gamma_n(\{x+t\})\dd x$ with the two ambient cutoffs.
For clarity, the matching to a finite part can be justified before
extracting any coefficients. Put
$\mathcal H(u,x)=\zeta(1+u,x)-1/u$ and $b=1-t$.
The following locally convergent expression is analytic for small $u,v$:
\begin{align*}
 &\int_0^b\left\{
 x^{-1-u}[\mathcal H(v,t+x)-\mathcal H(v,t)]
       +\mathcal H(u,1+x)\mathcal H(v,t+x)\right\}\dd x\\
 &+\int_0^t\left\{
 y^{-1-v}[\mathcal H(u,b+y)-\mathcal H(u,b)]
       +\mathcal H(v,1+y)\mathcal H(u,b+y)\right\}\dd y\\
 &+\mathcal H(v,t)\frac{1-b^{-u}}u
       +\mathcal H(u,b)\frac{1-t^{-v}}v.
\end{align*}
The bracketed differences vanish linearly; both quotients have removable
values at zero. The coefficient of $u^mv^n$ is exactly
$(-1)^{m+n}I_{mn}(t)/(m!n!)$, by direct integration of the subtracted
logarithmic singularities. In a left spectral half-plane the displayed
expression is also
\[
 C(1+u,1+v;t)+\frac1{uv}
       +\frac{\mathcal H(v,t)}u+\frac{\mathcal H(u,b)}v.
\]
Using the axis values of \eqref{mld:dil:K-hurwitz}, it equals
\[
 \frac{\mathcal K(u,v;t)-\mathcal K(u,0;t)
            -\mathcal K(0,v;t)+\mathcal K(0,0;t)}{uv}.
\]
This proves \eqref{mld:dil:I-coefficient} with the specified normalization.

The coefficient calculation uses rational polynomial arithmetic and
positive integer zeta values only, because
\begin{equation}
 \log\mathcal A(u,v)=\sum_{j\ge2}\frac{\zeta(j)}j
       \left\{(-1)^j\bigl((u+v)^j-u^j\bigr)+v^j\right\}.
 \label{mld:dil:logA}
\end{equation}
For an explicit finite algorithm, let $\ell_j$ be the homogeneous summand
of degree $j$ here. The homogeneous terms of $\mathcal A$ satisfy
\[
 \mathcal A_0=1,\qquad\mathcal A_1=0,\qquad
 \mathcal A_d=\frac1d\sum_{j=2}^d j\ell_j\mathcal A_{d-j}.
\]
Only $d\le m+n+2$ can contribute to \eqref{mld:dil:I-coefficient}. This proves
finite linear closure in the single values $\gamma_j(t)$ and
$\gamma_j(1-t)$, with $j\le m+n+1$. Their top coefficients are
$1/(m+1)$ and $1/(n+1)$; the coefficients at order $m+n$ vanish since
$\mathcal A_1=0$. These are the incoming closure theorem and its
missing-order rule, now available as an explicit input to the new
dilation calculation.


\subsection{Unequal integer dilations: the full Stieltjes theorem}
Fix positive integers $p,q$, put
\begin{equation}
 d=\gcd(p,q),\qquad P=p/d,\qquad Q=q/d,\qquad
 \theta=\{Pa\},\qquad 0<\theta<1.
 \label{mld:dil:parameters}
\end{equation}
The last condition is exactly the condition that the singular grids of
$\gamma_m(\{px\})$ and $\gamma_n(\{qx+a\})$ do not intersect.
Indeed an intersection would give $qx+a\in\mathbb Z$ and
$px\in\mathbb Z$, equivalently $Pa\in\mathbb Z$ since $(P,Q)=1$.
Let
\[
 F_{mn}^{p,q}(\theta)
 =\operatorname{FP}_x\int_0^1
       \gamma_m(\{px\})\gamma_n(\{qx+a\})\,dx.
\]
The notation $\operatorname{FP}_x$ means the common ambient-coordinate
constant term specified above. One may equivalently take a separate
constant term in an independent ambient cutoff variable at each grid
point. Replacing those variables by $c_j\epsilon$ and taking a single
constant term would generally change the answer through local
$\log c_j$ terms.

We use the already established undilated correlation
\[
 I_{mn}(t)=\operatorname{FP}\int_0^1
                  \gamma_m(x)\gamma_n(\{x+t\})\,dx,
 \qquad 0<t<1.
\]
Its all-order single-Stieltjes closure and finite coefficient rule were
recalled in Section~\ref{mld:dil:undilated-algorithm}.  The formulas needed below are
\begin{align}
 I_{00}(t)&=\gamma_1(t)+\gamma_1(1-t)-2\zeta(2),
 \label{mld:dil:I00}\\
 I_{10}(t)&=\tfrac12\gamma_2(t)+\gamma_2(1-t)
        +\zeta(2)\{\gamma_0(t)+\gamma_0(1-t)\}-\zeta(3).
 \label{mld:dil:I10}
\end{align}
For $M\ge1$ define a normalized trace polynomial
\begin{equation}
 A_n(M,t)=\sum_{j=0}^n\binom nj(-\log M)^{n-j}\gamma_j(t)
                 -\frac{(-\log M)^{n+1}}{n+1}.
 \label{mld:dil:trace-polynomial}
\end{equation}
In particular $A_0(M,t)=\gamma_0(t)+\log M$.
The pointwise Hurwitz multiplication identity \cite[\S25.11]{DLMFhurwitz} proves
\[
 \frac1M\sum_{j=0}^{M-1}\gamma_n\!\left(\left\{u+\frac jM\right\}\right)
                     =A_n(M,\{Mu\}),\qquad Mu\notin\mathbb Z.
\]

\begin{theorem}[Unequal-dilation Stieltjes closure]
\label{mld:dil:full-stieltjes}
Set $A=\log P$, $B=\log Q$, $L_p=\log p$, and $L_q=\log q$.
For all $m,n\ge0$ and the parameters \eqref{mld:dil:parameters},
\begin{align}
 F_{mn}^{p,q}(\theta)={}&
 \sum_{i=0}^m\sum_{j=0}^n\binom mi\binom nj
             (-B)^{m-i}(-A)^{n-j}I_{ij}(\theta)\notag\\
 &+\frac{(-B)^{m+1}-L_p^{m+1}}{m+1}A_n(P,\theta)
  +\frac{(-A)^{n+1}-L_q^{n+1}}{n+1}A_m(Q,1-\theta)\notag\\
 &+\frac{(-B)^{m+1}(-A)^{n+1}}{(m+1)(n+1)}.
 \label{mld:dil:full-formula}
\end{align}
Thus every such product is a finite linear combination of single
Stieltjes values at $\theta$ and $1-\theta$, with coefficients in the
algebra generated over $\mathbb Q$ by the logarithms of $p,q,P,Q$ and
positive integer zeta values.
\end{theorem}

\begin{proof}
Let $C(s,t;\theta)$ be the ordinary translated Hurwitz correlation,
initially for $\Re s,\Re t<1$.  Fourier orthogonality first in
$\Re s,\Re t<0$ gives
\begin{equation}
 \int_0^1\zeta(s,\{px\})\zeta(t,\{qx+a\})\,dx
                   =Q^{s-1}P^{t-1}C(s,t;\theta).
 \label{mld:dil:frequency}
\end{equation}
Indeed the surviving frequency indices satisfy $pu+qv=0$, hence
$u=Q\ell$, $v=-P\ell$; their translation factor is
$e^{-2\pi i\ell Pa}$.  The two sides have the same meromorphic
continuation after local endpoint subtraction.  The common divisor $d$
does not occur in \eqref{mld:dil:frequency}, because ordinary integration is
unchanged by a common covering.

Put $\mathcal G(z,x)=\sum_{n\ge0}\gamma_n(x)z^n/n!$ and
$\mathcal I(z,w;\theta)=\sum_{m,n\ge0}I_{mn}(\theta)z^mw^n/(m!n!)$.
For distinct singular points the distributional Laurent series gives
\begin{equation}
 C(1-z,1-w;\theta)
 =-\frac1{zw}+\frac{\mathcal G(w,\theta)}z
                +\frac{\mathcal G(z,1-\theta)}w
                +\mathcal I(z,w;\theta).
 \label{mld:dil:C-Laurent}
\end{equation}
For example, the coefficient $-1/(zw)$ is the correlation of
$\delta_0-1$ with its nontrivial translate.  Omitting it would change
the final constant in \eqref{mld:dil:full-formula}.

Multiply \eqref{mld:dil:C-Laurent} by $e^{-Bz-Aw}$ and extract
$z^mw^n/(m!n!)$.  This computes the correlation of the \emph{true
pullbacks} $U_pG_m$ and $G_n(qx+a)$.  Its value is the double sum in
\eqref{mld:dil:full-formula}, plus
\[
 \frac{(-B)^{m+1}}{m+1}A_n(P,\theta)
 +\frac{(-A)^{n+1}}{n+1}A_m(Q,1-\theta)
 +\frac{(-B)^{m+1}(-A)^{n+1}}{(m+1)(n+1)}.
\]
The $r=0$ dilation law changes the first factor by
$-L_p^{m+1}\Delta_p/(p(m+1))$.  Pairing this with the other factor
gives $-L_p^{m+1}A_n(P,\theta)/(m+1)$: the grid values form $d$
copies of a complete reduced $P$-grid.  The second factor contributes
$-L_q^{n+1}A_m(Q,1-\theta)/(n+1)$ similarly.  The product of the two
grid distributions vanishes, because their supports are disjoint.
Their sum is \eqref{mld:dil:full-formula}.  The product and all coefficient
extractions can be justified using a partition of unity separating the
two singular grids: on each singular neighborhood the other factor is
smooth, and elsewhere both factors are smooth.
\end{proof}

For example, if $p=2$, $q=1$, $\theta=\{2a\}\in(0,1)$ and
$L=\log2$, \eqref{mld:dil:full-formula} gives the explicit new family member
\begin{align}
 \operatorname{FP}_x\int_0^1\gamma_1(\{2x\})\gamma_0(\{x+a\})\,dx
 ={}&\tfrac12\gamma_2(\theta)+\gamma_2(1-\theta)
        -L\gamma_1(1-\theta)\notag\\
 &+\left(\zeta(2)-\tfrac12L^2\right)\gamma_0(\theta)
        +\zeta(2)\gamma_0(1-\theta)-\zeta(3)-\tfrac12L^3.
 \label{mld:dil:mixed-example}
\end{align}

\subsection{Every mixed parameter derivative}
\begin{theorem}[Dilated differentiation and product integration]
\label{mld:dil:all-derivatives}
Let $N=r+k$ with $r,k,m,n\ge0$, and differentiate
$F_{mn}^{p,q}$ in its displayed argument $\theta$.  Then
\begin{align}
 &\operatorname{FP}_x\int_0^1
       \gamma_m^{(r)}(\{px\})\gamma_n^{(k)}(\{qx+a\})\,dx\notag\\
 &=P^kQ^r\biggl\{
       (-1)^r\frac{d^N}{d\theta^N}F_{mn}^{p,q}(\theta)
       +(-1)^r b_{r,m}(L_p)\partial_\theta^N A_n(P,\theta)\notag\\
 &\hspace{40mm}
       +(-1)^k b_{k,n}(L_q)
          \left.\partial_t^N A_m(Q,t)\right|_{t=1-\theta}
               \biggr\}.
 \label{mld:dil:all-derivative-formula}
\end{align}
The derivatives on the left are derivatives in the \emph{special-function
argument}, not derivatives of the fractional-part map.
\end{theorem}
\begin{proof}
Rewrite \eqref{mld:dil:commutator} as
\[
 X_{p,m,r}=p^{-r}D^rX_{p,m,0}
          +p^{-r-1}b_{r,m}(L_p)\Delta_p^{(r)},
\]
and use the translated counterpart for the second factor.  Periodic
integration by parts in the first product yields
\[
 p^{-r}q^{-k}\int(D^rX_{p,m,0})(D^kX_{q,n,0}(\,\cdot+a/q))
 =(-1)^r(q/p)^r\partial_a^NF_{mn}^{p,q}(\{Pa\}).
\]
Since $\partial_a=P\partial_\theta$ on each noncollision interval,
this is $(-1)^rP^kQ^r\partial_\theta^NF$.
The first grid contact pairs with the smooth second factor and gives
\[
 (-1)^r(q/p)^r b_{r,m}(L_p)\frac1p
       \sum_{j=0}^{p-1}\gamma_n^{(N)}(\{qj/p+a\})
 =(-1)^rP^kQ^r b_{r,m}(L_p)\partial_\theta^NA_n(P,\theta).
\]
The second contact is the third term in
\eqref{mld:dil:all-derivative-formula}.  The grid-grid term vanishes by
disjointness.  The same separated partition of unity used in the
preceding proof makes these integrations by parts valid.
\end{proof}

\subsection{A particularly short polygamma identity}
Put $H_0=0$, $H_j=\sum_{\ell=1}^j1/\ell$ for $j\ge1$, and
\[
 K=\log\operatorname{lcm}(p,q).
\]
\begin{corollary}[All unequal-dilation polygamma correlations]
\label{mld:dil:polygamma}
For $N=r+k\ge1$,
\begin{align}
 &\operatorname{FP}_x\int_0^1
         \psi^{(r)}(\{px\})\psi^{(k)}(\{qx+a\})\,dx\notag\\
 &=N!P^kQ^r\biggl\{
 (-1)^{k+1}\bigl[\zeta'(N+1,\theta)
                 +(H_N-H_r+K)\zeta(N+1,\theta)\bigr]\notag\\
 &\hspace{24mm}+(-1)^{r+1}\bigl[\zeta'(N+1,1-\theta)
                 +(H_N-H_k+K)\zeta(N+1,1-\theta)\bigr]\biggr\}.
 \label{mld:dil:polygamma-formula}
\end{align}
The two primes on the right denote first derivatives in the spectral
variable of Hurwitz zeta.  The $r=k=0$ case is
\begin{align}
 &\operatorname{FP}_x\int_0^1\psi(\{px\})\psi(\{qx+a\})\,dx\notag\\
 &=\gamma_1(\theta)+\gamma_1(1-\theta)
       +K\{\psi(\theta)+\psi(1-\theta)\}
       -\frac{\pi^2}{3}-K^2+\log p\log q.
 \label{mld:dil:digamma-formula}
\end{align}
\end{corollary}
\begin{proof}
Apply \eqref{mld:dil:full-formula} with $m=n=0$.  Its logarithmic
coefficient is $K=L_p+B=L_q+A$, and its constant simplifies using
$AB-K(A+B)=-K^2+L_pL_q$.  This proves
\eqref{mld:dil:digamma-formula}.  In
\eqref{mld:dil:all-derivative-formula}, use $b_{r,0}=H_r$,
$\gamma_0=-\psi$, and the pointwise identities
\[
 \psi^{(N)}(t)=(-1)^{N+1}N!\zeta(N+1,t),\qquad
 \gamma_1^{(N)}(t)=(-1)^{N+1}N!
       \{\zeta'(N+1,t)+H_N\zeta(N+1,t)\}.
\]
The two reflection signs are $(-1)^r$ and $(-1)^k$ before this
conversion.  Simplification gives \eqref{mld:dil:polygamma-formula}.
\end{proof}

The least common multiple in the logarithm is forced by the local
normalization.  For example, a common dilation $p=q=d$ leaves the
ordinary spectral correlation unchanged, but its digamma finite part
changes by $\log d\{\psi(\theta)+\psi(1-\theta)\}$.
Discarding this term would confuse transported finite parts with
ambient-coordinate finite parts.

\begin{corollary}[Odd total derivative order at a half shift]
If $\{Pa\}=1/2$ and $N=r+k$ is odd, then
\begin{align}
 &\operatorname{FP}_x\int_0^1
       \psi^{(r)}(\{px\})\psi^{(k)}(\{qx+a\})\,dx\notag\\
 &=N!P^kQ^r(-1)^{k+1}(H_k-H_r)
                       (2^{N+1}-1)\zeta(N+1).
 \label{mld:dil:half-odd}
\end{align}
In particular every such value is a rational multiple of
$\pi^{N+1}$.  One unequal-dilation example is
\begin{equation}
 \operatorname{FP}_x\int_0^1
       \psi(\{2x\})\psi'(\{3x+1/4\})\,dx=\pi^2.
 \label{mld:dil:pi2-example}
\end{equation}
\end{corollary}
\begin{proof}
When $N$ is odd the two signs in \eqref{mld:dil:polygamma-formula} are
opposite.  At $\theta=1/2$ the spectral derivatives, the $H_N$ terms,
and the $K$ terms cancel pairwise.  The surviving harmonic difference
multiplies $\zeta(N+1,1/2)=(2^{N+1}-1)\zeta(N+1)$.
Since $N+1$ is even, Euler's even-zeta formula gives the last assertion.
\end{proof}

The same unequal pair at zero derivative order gives a different exact
identity, now involving the ordinary first Stieltjes constant:
\begin{align}
 &\operatorname{FP}_x\int_0^1
       \psi(\{2x\})\psi(\{3x+1/4\})\,dx\notag\\
 &=2\gamma_1-2\gamma\log24-\frac{\pi^2}{3}
       -7\log^22-5\log2\log3-\log^23.
 \label{mld:dil:half-digamma-example}
\end{align}
This follows from \eqref{mld:dil:digamma-formula},
$\psi(1/2)=-\gamma-2\log2$ and
$\gamma_1(1/2)=\gamma_1-2\gamma\log2-\log^22$.

\subsection{Ordinary convergent-integral realization}
Every identity in Corollary~\ref{mld:dil:polygamma} can be written without
a finite-part symbol.  Sort the two disjoint singular grids as
$0=x_0<x_1<\cdots<x_{M-1}<1$, put $x_M=1$, and set
$\ell_i=x_{i+1}-x_i$.  On each interval use the ordinary branches
\[
 f_i(t)=\psi^{(r)}(\alpha_i+pt)
        \psi^{(k)}(\beta_i+qt),\qquad0<t<\ell_i,
\]
where $\alpha_i=\{px_i\}$ and $\beta_i=\{qx_i+a\}$, with the
right-hand branch chosen at zero.  Exactly one of $\alpha_i,\beta_i$
is zero.  If $\alpha_i=0$, set $v_i=r$ and
\[
 a_{ij}=\frac{(-1)^{r+1}r!q^j}{p^{r+1}j!}
                       \psi^{(k+j)}(\beta_i),\quad0\le j\le r.
\]
If $\beta_i=0$, set $v_i=k$ and
\[
 a_{ij}=\frac{(-1)^{k+1}k!p^j}{q^{k+1}j!}
                       \psi^{(r+j)}(\alpha_i),\quad0\le j\le k.
\]
Then the left side of \eqref{mld:dil:polygamma-formula}, including its
$N=0$ counterpart, equals the ordinary absolutely convergent expression
\begin{equation}
 \sum_{i=0}^{M-1}\left[
  \int_0^{\ell_i}\left\{f_i(t)
               -\sum_{j=0}^{v_i}a_{ij}t^{j-v_i-1}\right\}\,dt
  +a_{i,v_i}\log\ell_i
  +\sum_{j=0}^{v_i-1}\frac{a_{ij}\ell_i^{j-v_i}}{j-v_i}
                           \right].
 \label{mld:dil:ordinary-realization}
\end{equation}
This follows by integrating the displayed singular monomials explicitly.
The subtracted integrands are bounded at their left endpoints.  At the
right endpoint the branch approaching one is used, so no singularity
is introduced there.

The independent numerical implementation uses precisely
\eqref{mld:dil:ordinary-realization} on the left and
\eqref{mld:dil:polygamma-formula}--\eqref{mld:dil:digamma-formula} on the right.
Nine cases include unequal dilations, a nontrivial common divisor, and
derivative orders up to $r+k=3$.  At $42$ decimal working digits the
largest relative discrepancy is below $3.5\times10^{-39}$.  These are
floating-point consistency checks of formulas proved above, not certified
error enclosures or a substitute for the analytic argument.
Three further independent subtracted-integral checks use positive
Stieltjes indices $(m,n)=(1,0),(0,1),(1,1)$, the last with
$(p,q)=(2,3)$.  At $28$ working digits their largest relative discrepancy
is below $9\times10^{-29}$.


% SOURCE: sections/06-audit-research.tex
\section{Research status, corrections, and further questions}
\label{mld:sec:research}

\subsection{What the continuation settles}

The main resolved question in the supplied incoming set concerns integer
dilation and the choice of local subtraction coordinate. Theorem~
\ref{mld:dil:scale-thm} gives the complete local correction, including every
Stieltjes index and argument derivative. Equations~\eqref{mld:dil:full-formula}
and \eqref{mld:dil:all-derivative-formula} then evaluate unequal-dilation
products when the singular grids are disjoint. Their dependence on
$\log\lcm(p,q)$ is explicit. These statements answer the integer-covering
case of the rescaling question; a general nonlinear coordinate change
is outside their scope.

The exact locations are the further-research items on unequal dilations
in \emph{Bilinear Closure of Generalized Stieltjes Correlations}
\cite{IncomingClosure}, on dilation, traces and subtraction scale in
\emph{Shifted Hurwitz-Zeta Correlations} \cite{IncomingShifted}, and on
rescaling in \emph{A Convolution Calculus for Stieltjes Functions}
\cite{IncomingCalculus}. The integration notes give the archive members
and line locations in the pinned snapshot.

The Mellin results are independent of that regularization problem.
They give ordinary integrals with a uniform finite formula at arbitrary
denominator order. The integer resonances, all their logarithmic moments,
and the spectral derivatives in Section~\ref{mld:sec:jets} are proved
identities. The harmonic reflection is likewise an identity of
holomorphic functions and ordinary convergent series. Its even-weight
unshifted coefficient evaluations have published antecedents
\cite{XuWang}; the two-shift formulation is presented with a proof and
without an unverified claim of historical priority.

\begin{center}
\begin{tabularx}{\textwidth}{@{}>{\raggedright\arraybackslash}p{0.35\textwidth}Y@{}}
\toprule
Claim & Status after this work\\
\midrule
All-$N$ Mellin transform and integer resonances & Proved in
Theorems~\ref{mld:thm:Mellin-master} and \ref{mld:thm:integer-resonance}.\\
Shifted generalized-harmonic reflection & Proved in
Theorem~\ref{mld:gh:thm:reflection}; scalar parity precedent credited.\\
Integer-dilation finite parts and disjoint-grid products & Proved in
Theorems~\ref{mld:dil:scale-thm}, \ref{mld:dil:full-stieltjes}, and
\ref{mld:dil:all-derivatives}.\\
Undilated Stieltjes closure and convolution & Already proved in the
supplied incoming material; recalled as inputs.\\
Gaussian $S_4$ evaluation & Already proved in the pinned manuscript;
its status is unchanged.\\
Current Gaussian $S_6$ and $S_8$ evaluations & Still conjectural.
No equality follows from a small residual.\\
Colliding unequal singular grids; nonlinear coordinate changes &
Not treated by the new dilation theorem.\\
\bottomrule
\end{tabularx}
\end{center}

\subsection{The Gaussian conjectures and the missing parity component}

The manuscript defines
\[
 S_p=\sum_{n\ge0}\frac{(-1)^nH_n}{(2n+1)^p},\qquad
 g_{a,b}=\operatorname{Im}\Li_{a,b}(\ii,1),
\]
where
$\Li_{a,b}(x,y)=\sum_{n>m\ge1}x^ny^m/(n^am^b)$ in its convergence
region. Its proved translation is
\[
 S_p=\operatorname{Im}\bigl(\Li_{p,1}(\ii,1)
                                  +\Li_{p,1}(\ii,-1)\bigr).
\]
The file \texttt{04-S4-proof.tex} supplies the exact $S_4$ certificate
and a convergent octahedral relation. The current $S_6$ conjecture is in
\texttt{04-cyclotomic-quotients.tex}; the current $S_8$ conjecture is in
\texttt{04-S8-candidate.tex}. Their complete coefficient vectors and
source labels are preserved in \texttt{integration/GAUSSIAN\_STATUS.md}.

For specificity, the $S_6$ equality still to be proved is
\begin{align}
 S_6\stackrel{?}{={}}&\frac{722}{527}g_{6,1}
  +\frac{40}{527}g_{4,3}-\frac{128}{155}g_{2,5}
  +\frac{15191}{28569600}\pi^7\notag\\
 &-\frac3{124}\beta(2)\zeta(5)
  -\frac{2373}{10540}\beta(4)\zeta(3)-2\beta(6)\log2.
 \label{mld:eq:S6-open}
\end{align}
The manuscript also proves a change of coordinates between two
presentations of this candidate. That identity equates two residual
expressions; it does not show that either residual is zero.

The obstruction to proving \eqref{mld:eq:S6-open} by the reflection obtained
here is exact. Its Taylor coefficient is
\[
 [z^{p-1}u^{r-1}]\{A(z,u)+A(-z,-u)\}
           =(1+(-1)^{p+r})Q_{p,r}.
\]
Since $S_6=Q_{6,1}$ and $S_8=Q_{8,1}$ have total weights $7$ and $9$,
their coefficients vanish in this projection. The same is true of the
weight-$5$ $S_4$ value: its known proof uses additional functional
information. Further differentiation of the same even projection
does not recover its absent Taylor coefficients. A successful extension
must control the complementary function
$A(z,u)-A(-z,-u)$ or produce an independent relation among the required
mixed-color periods.

The separating functional already exhibited in the manuscript rules out
a particular restricted row-span certificate for the $S_6$ candidate.
It does not rule out the identity itself. Similarly, a formal quotient
dimension does not prove arithmetic independence of the numerical
periods. These distinctions are retained throughout this continuation.

\subsection{Audit findings and concrete integration corrections}

No new false mathematical statement was found in the sampled pinned
Gaussian claims or the incoming correlation identities used here.
Several concrete reconciliations are nevertheless needed when the
packages are integrated.

\begin{enumerate}[label=\arabic*.]
\item \textbf{Reconcile the open-question lists.}
The Closure report's questions about undilated parameter derivatives and
associative convolution are answered in the current Shifted Hurwitz and
Convolution packages. Their entries should point to those results.
The integer-dilation questions should point to the present scale and
product theorems, with disjointness and ambient normalization stated.

\item \textbf{Keep the two dilation operations distinct.}
The replacement $X_{q,n,r}=U_qT_{n,r}$ would omit
\eqref{mld:dil:scale-law}. A simple exact diagnostic is
\[
 \langle U_qG_0,1\rangle=0,\qquad
 \langle X_{q,0,0},1\rangle=-\log q.
\]
Thus they already differ at zero Stieltjes index and zero derivative
order. This corrects a possible naive extension of the incoming formulas;
it does not assert that the supplied undilated theorem made that error.

\item \textbf{State the cutoff variable in derivative identities.}
The valid undilated assertion $h_{r,n}=0$ for $n\ge r$ cannot be copied
unchanged into a theorem about a fixed ambient scale. For example,
\[
 DX_{q,1,0}=qX_{q,1,1}-\frac{\log q}{q}\Delta_q'.
\]
A separate constant term in each local variable has the canonical
normalization used here. Rescaling those variables before one common
constant-term extraction changes it.

\item \textbf{Retain resonance values before simplifying factors.}
At integer spectral parameters, use the analytic germ
\eqref{mld:eq:cancelled-zeta-germ}. The equality $F_0(1)=1$ is responsible
for the $-1/4$ and $-3/4$ higher-pole examples. Dropping a factor that
vanishes next to a zeta pole is not a valid simplification.

\item \textbf{Preserve the existing conjecture and error history.}
The present $S_8$ candidate is different from the earlier vector already
rejected in \texttt{10-discovery.tex}. The mixed-color symmetry
counterexample is already recorded in \texttt{04-depth.tex}.
These are existing corrections, not new discoveries of this audit.
The generalized-harmonic family $Q_{p,r}$ should also retain its own
notation, since the manuscript's $S_{p,r}$ uses elementary symmetric
harmonic sums.
\end{enumerate}

\subsection{Proposed research programme}

The following questions seek exact identities or structural reductions.
They are research questions, not conjectures inferred from a few
numerical samples. The first three are the most immediate next steps.

\begin{enumerate}[label=\textbf{Q\arabic*.},leftmargin=2.7em]
\item \textbf{The missing Gaussian projection.}
Find a functional equation or convergent iterated-integral relation that
determines the weight-$7$ and weight-$9$ coefficients of
$A(z,u)-A(-z,-u)$. A concrete success criterion is an exact zero
certificate for \eqref{mld:eq:S6-open}, followed by the current $S_8$ vector.
The known $S_4$ octahedral seed suggests testing genuinely new functional
relations before enlarging a numerically fitted coefficient basket.

\item \textbf{Colliding unequal dilations.}
Extend the ambient-coordinate product theorem to $Pa\in\Z$. There the
partition separating the two singular grids fails, and the product
contains higher singular powers at the common points. Derive the full
local subtraction polynomial and compare its finite value with the
collision expansion of \eqref{mld:dil:full-formula}. The already known
undilated coincident closure is a boundary case and a normalization
check, not a solution of every unequal-grid configuration.

\item \textbf{Nonlinear changes of local coordinate.}
Replace $qt$ near a singularity by an analytic map
$f(t)=qt+c_2t^2+\cdots$, $q>0$. Determine the finite combination of
delta derivatives that converts pullback into ambient finite part for
every $\gamma_n^{(r)}$. The linear coefficient must recover
$c_{n,r}(\log q)$; higher coefficients should be computable from finite
Taylor and Bell-polynomial algebra. The goal is a proved composition
law for the full local correction.

\item \textbf{Three translated or dilated factors.}
For three pairwise disjoint singular grids, determine the spectral
generating kernel of the product of three Hurwitz or Stieltjes
functions. Fourier orthogonality leaves a two-dimensional frequency
constraint. Identify which double Hurwitz sums or multiple
polylogarithms survive, and derive the first nontrivial exact Stieltjes
coefficient formulas. Do not assume the bilinear single-Stieltjes
closure persists at degree three.

\item \textbf{Complete primitive calculus for dilated products.}
Insert the finite closure into the anchored Hurwitz-jet primitive
ladder to construct explicit primitives in the shift $a$ on every
noncollision interval. Then determine the matching constants when
crossing a collision point in a specified finite-part convention.
This would connect local exact integration with a global periodic
distribution statement.

\item \textbf{Character-weighted dilation and multiplication.}
Replace the unweighted trace by Dirichlet-character weights on the
finite covering. Determine how the local correction interacts with
the character's zero mode, and express the resulting correlations in
Dirichlet $L$-jets. The residue-class trace polynomial
\eqref{mld:dil:trace-polynomial} supplies the unweighted test case.

\item \textbf{Rational-shift harmonic identities beyond a reflection.}
Combine the two-shift theorem with recurrences in its inner and outer
shifts. Determine which finite combinations of individual shifted
harmonic sums admit depth-one cyclotomic expressions, rather than only
their reflected pairs. A useful first classification fixes small
denominators in $u$ and $z$ and records exact functional certificates.

\item \textbf{Mellin kernels with noninteger denominator order.}
Replace $N$ in $(1+x)^{-N}$ by a complex parameter. Find a spectral
representation that specializes to the finite polynomial operator at
positive integers, and determine exactly when a finite ordinary-
polylogarithm closure survives. Hypergeometric or generalized Lerch
functions are natural intermediate objects, but their branch and
resonance conventions must be derived from the convergent integral.

\item \textbf{Complex poles and contour rotation.}
Extend Corollary~\ref{mld:cor:rational-kernel} to rational kernels with
nonreal poles. Track the residue terms when a rotated integration ray
crosses a polylogarithm cut. The target is an exact continuation formula
that reduces to the positive-real scaling law when every pole lies on
the negative real axis.

\item \textbf{Products of polylogarithms under the Mellin kernel.}
Study
\[
 \int_0^\infty \frac{x^{a-1}\Li_s(-zx)\Li_t(-wx)}{(1+x)^N}\dd x
\]
where ordinary convergence holds. Decide whether integer Mellin
parameters lead to a finite multiple-polylogarithm algebra, and identify
the minimal depth forced by two independent arguments. The one-factor
integer table gives exact boundary cases and a recurrence in $N$.

\item \textbf{Compatibility of the two harmonic coefficient calculi.}
Both cancelled zeta poles and dilation corrections involve
$\prod_j(1-u/j)$, whereas repeated anchored primitives involve its
reciprocal. Formulate and prove a finite algebra connecting these
operations, including derivative--primitive commutators and their
integration polynomials. This could reduce duplicated coefficient
calculations across the manuscript.

\item \textbf{Exact certificates and formal verification.}
Encode the polynomial recurrences, the scale composition law, and the
integer moment table as exact certificates with stable normalization
metadata. A proof-assistant development could then isolate the finite
algebra from the analytic tasks of dominated parameter differentiation,
meromorphic continuation, and distributional pullback. Numerical
agreement alone should not change theorem status.
\end{enumerate}

These questions are constrained by the proved formulas: each has a
specified base case, normalization, or missing component to resolve.
They also preserve the main distinction of this article: ordinary
polylogarithmic integrals, finite harmonic reductions, and normalized
distribution products are connected by exact operations, but are not
interchangeable without their hypotheses.
